\pagebreak

\chapter{Variables, Constants, Types, Modes, Attributes, Operators and Expressions}



    \section{Input format}
    
        \index{input format}
        The source code input may use either UNIX-style (newline character) or
        DOS/Windows-style (carriage return, line feed) line termination. Tokens
        are separated by white space (spaces, horizontal tabs or newlines). Good
        formatting dictates the use of indenting but the compiler treats
        indenting as simply white space. The last source code line need not be
        terminated by a newline character, but it is usually preferable to do
        so.

        \index{comments}
        Comments may have the form "/* ....... */", in which case they
        may run over line breaks, or they may be indicated by "//....."
        where they occupy the rest of the current line.


    \section{File inclusion}\label{sec:fileinclusion}

        Source files may be inserted in-line by means of file inclusion
        statements. These are of two forms.
        
        The first form
        \begin{lstlisting}[gobble=8, language=threepl]
           source "filename"
        \end{lstlisting}
        \index{source statement}
        \index{statement!source}
        simply inserts the contents of the named file into the current file in
        place of the entire {\em source} statement line. A file may only be
        inserted once by a {\em source} statement.
        
        The second form
        \begin{lstlisting}[gobble=8, language=threepl]
           module "filename"
        \end{lstlisting}
        \index{module statement}
        \index{statement!module}
        also inserts the contents of the named file into the current file in
        place of the entire {\em module} statement line, but it differs from a
        {\em source} statement in that the inserted file establishes a new
        static module scope (for scope discussion see "Program source structure"
        \autorefph{chap:progsrcstruct}). A file may only be inserted once by a
        {\em module} statement.

        White space (only) may precede the {\em source} or {\em module} on the
        line and separate it from the quoted file name. The remainder of the line
        following the quoted file name will be ignored. File inclusions using {\em
        source} or {\em module} statements may be nested to any depth.

        The quoted file name may be replaced in either of the above by
        {\bf directive("ident")}, i.e.
        \begin{lstlisting}[gobble=8, language=threepl]
           source directive("key")
        \end{lstlisting}
        or 
        \begin{lstlisting}[gobble=8, language=threepl]
           module directive("key")
        \end{lstlisting}
        where the directive whose key is "key" contains the name of the
        file to be included. Since file inclusion is done during parsing
        the directive must be assigned using {\bf presetdirective()}
        \autorefph{sec:ippresetdirective} or by using the command line option
        -D.
        
        \index{ifdirective}
        The file inclusion may be made conditional by using one of the following
        forms -
        \begin{lstlisting}[gobble=8, language=threepl]
           source ifdirective("key") "ident"
        \end{lstlisting}
        or 
        \begin{lstlisting}[gobble=8, language=threepl]
           module ifdirective("key") "ident"
        \end{lstlisting}
        where the directive "key" will be examined to determine whether to
        include the file "ident". If directive "key" does not exist or
        it has the boolean value {\bf false} the file inclusion will
        be skipped, otherwise the file will be included. Note that if the
        directive is not boolean its existence will suffice, but if
        boolean its value will be used.
        
        The above forms may also use a directive for the file identifier -
        \begin{lstlisting}[gobble=8, language=threepl]
           source ifdirective("key1") directive("key2")
        \end{lstlisting}
        or 
        \begin{lstlisting}[gobble=8, language=threepl]
           module ifdirective("key1") directive("key2")
        \end{lstlisting}
        Here directive "key1" is used as before to determine whether to
        include the file and directive "key2" contains the name of the
        file. The same directive may be used for both purposes, e.g.
        \begin{lstlisting}[gobble=8, language=threepl]
           module ifdirective("lib1") directive("lib1")
        \end{lstlisting}
        will include the file whose name is contained in directive "lib1"
        if that directive is defined.

        File names are used exactly as supplied so any file name extension
        must be included (use of the extension ".3pl" is recommended but not
        required).

        Note that these file inclusions are performed during parsing, not during
        immediate execution, and so placing a {\em source} statement within
        an {\em if} statement does not make the inclusion conditional (and
        is syntactically incorrect).
        
        Files included by either of these mechanisms must contain syntactically
        complete code, i.e. statements cannot be partially contained in an
        included file.
        
        Since there are two mechanisms for inserting files, a file should
        be able to ensure that it is not inserted by the wrong mechanism.
        This can be done by including one of the lines
        \begin{lstlisting}[gobble=8, language=threepl]
           nosource
        \end{lstlisting}
        \index{nosource statement}
        \index{statement!nosource}
        to preclude sourcing this file or
        \begin{lstlisting}[gobble=8, language=threepl]
           nomodule
        \end{lstlisting}
        \index{nomodule statement}
        \index{statement!nomodule}
        to preclude calling this file as a module.
        These must have no leading characters before the keyword. The remainder
        of the line is ignored. 



    \section{Single constants}

        \index{constants}
        Integer constants may be
        \blist
        \item    decimal
        \item    hexadecimal (leading 0x or 0X)
        \item    octal (leading 0)
        \item    binary (leading 0b or 0B).
        \blend
        Their size is 64 bits.

        Logical constants are {\bf true} and {\bf false}.

        Floating point constants are decimal integers with a leading,
        trailing or embedded '.' and an optional exponent consisting of "E"
        or "e", an optional sign followed by an integer. If an exponent is
        supplied the '.' can be omitted where appropriate.

        String constants have double quotes surrounding the string. The string
        may contain escape characters
        \blist
        %\item   \verb'\a' alert (bell)
        %\item   \verb'\e' escape
	\item   \verb'\n' new line (line feed)
        \item   \verb'\t' horizontal tab
        %\item   \verb'\v' vertical tab
        \item   \verb'\b' backspace
        \item   \verb'\r' carriage return
        \item   \verb'\f' form feed
        \item   \verb'\\' backslash
	\item   \verb+\'+ single quote (single quotes do not need to be escaped)
        \item   \verb'\"' double quote
        \item   \verb'\0nnn' octal numeric escape where nnn is the 3 digit octal
                character code
        \item   \verb'\xnn' hexadecimal numeric escape where nn is the 2 digit
                hexadecimal character code
        \blend
        
        \index{mathematical constants}
        The global variables {\bf \verb'math_pi'} and {\bf \verb'math_e'}
        of type {\bf float} are created by the compiler prior to immediate
        execution. These variables contain the values for $\pi$ and $e$ and
        are read-only.

    \section{Compound values}

        \index{compound!value}
        \index{value!compound}
        Compound values are (possibly nested) array or struct values.
        \blist
        \item   They can be used to initialise compound variables.
        \item   They can be assigned to compound variables.
        \item   They can be passed to compound immediate or value
                parameters in modules, procedures or functions.
        \blend

        A compound value is a comma-separated list enclosed in braces. Each list
        element is an immediate expression or a compound value Each set of braces
        matches a top level or nested array or struct. Trailing elements for an
        array or struct may be omitted. Initial or intermediate elements may also
        be omitted, commas being placed to indicate this. Lists must not exceed the
        length of the array or struct they are intended to match.
        
        A compound value is processed from left to right, thus the left-most
        item will appear at the lowest subscript or address.
        
        Missing elements in a compound value are simply not assigned.
        Where a compound value with missing elements is used to initialise
        an immediate compound variable, or is assigned to it, the variable
        elements corresponding to the missing values will simply retain
        their previous or default values. Where a compound value is
        similarly used for a value variable, the variable elements
        corresponding to the missing constants will not be assigned, retaining
        their previous values (none if not previously initialised or assigned).

        Examples -
        \begin{lstlisting}[gobble=8, language=threepl]
           var(i, "[3][2]in", {{1, 2}, , {3}});
           struct(s,
               f1,     "[3]int",
               f2,     "log",
               f3,     "str"
           );
           var(v, s);
           v = {{3, 4, 5}, true, "fred"};
           proc(a <- {7, 8, false});
        \end{lstlisting}
        
        Where a struct is initialised, assigned or passed an argument the
        field values may be specified by naming the fields rather than matching
        values to fields by position as in the example above, e.g.
        \begin{lstlisting}[gobble=8, language=threepl]
           struct(s,
               f1,     "[3]int",
               f2,     "log",
               f3,     "str"
           );
           var(v, s);
           v = {f3="fred", f1={3, 4, 5}, f2=true};
           proc(a <- {7, 8, flag=false});
        \end{lstlisting}
        When a struct is assigned values by name in this way no following
        position-matched values are allowed, for example the following
        is accepted -
        \begin{lstlisting}[gobble=8, language=threepl]
           struct(s,
               f1,     "[3]int",
               f2,     "log",
               f3,     "str"
           );
           var(v, s);
           v = {{3, 4, 5}, f3="fred", f2=true};
        \end{lstlisting}
        but it is incorrect to do -
        \begin{lstlisting}[gobble=8, language=threepl]
           struct(s,
               f1,     "[3]int",
               f2,     "log",
               f3,     "str"
           );
           var(v, s);
           v = {{3, 4, 5}, f2=true, "fred"};
        \end{lstlisting}
        because the last value is matched by position but is preceded by
        a value matched by name. This limitation makes sense in that values
        are matched by position in the list to fields in the struct until
        a match by name is encountered after which point the argument positions
        are no longer aligned with the fields.
        
        \index{array!initialisation}
        \index{array!undimensioned}
        \index{dimensions!missing}
        An array initialised by a compound value may have some dimensions
        omitted, in which case the missing dimensions will be inferred from the
        compound value. Array dimensions cannot be omitted where a variable is
        not a module, procedure or function parameter and is not initialised.
        
        Compound value entries for a string array may be types integer,
        boolean or floating point in which case they will be converted to strings.

        Leaf elements in the compound value must be immediate primitive
        types (which includes pointers).


    \section{Identifiers}

        \index{identifier}
        Identifiers begin with a letter or an underscore and contain any
        number of letters, digits or underscores. Alphabetic case is
        distinguished.

        The following keywords may not be used as identifiers

        \index{keywords}
        \blist
        \item   branch
        \item   break
        \item   case
        \item   class
        \item   conditionelse
        \item   continue
        \item   default
        \item   do
        \item   else
        \item   for
        \item   function
        \item   global
        \item   if
        \item   ifcondition
        \item   leadermodule
        \item   module
        \item   null
        \item   par
        \item   petrinet
        \item   postmodule
        \item   private
        \item   procedure
        \item   return
        \item   seq
        \item   source
        \item   switch
        \item   sync
        \item   trailermodule
        \item   unless
        \item   varargs
        \item   when
        \item   while
        \blend

        Identifiers used for -
        \index{identifier!ambiguity}
        \blist
        \item   variables
        \item   struct fields
        \item   modules and procedures
        \item   functions
        \item   directives
        \blend
        are stored in separate tables and hence the same
        identifier can be used for multiple purposes. However, this practice may
        render the code difficult to read.








    \section{Variables and modes}

        \subsection{variable modes}

            \index{modes}
            There are eleven modes of variable -
            \blist
            \item    {\bf immediate}
            \item    {\bf selectvalue}
            \item    {\bf value}
            \item    {\bf static}
            \item    {\bf queue}
            \item    {\bf priority}
            \item    {\bf input}
            \item    {\bf output}
            \item    {\bf cmemory}
            \item    {\bf rmemory}
            \item    {\bf clock}
            \blend

            Immediate variables only exist during program interpretation.
            Variables in the other ten modes are target variables which
            are active in the FPGA logic.


            \subsubsection{immediate variables}

                \index{immediate!variable}
                \index{variable!immediate}
                \index{mode!immediate}
                Immediate variables do not survive into target code except as
                constants (their current value at that point in the code during
                immediate execution). These variables must not have a width
                declared. Immediate integer variables are 64 bits and floating
                point variables are 64 bit double precision.

            \subsubsection{value variables}

                \index{value!variable}
                \index{variable!value}
                \index{mode!value}
                \index{value mode variables}
                Variables of mode {\bf value} represent the result of
                combinatorial expressions and hence do not have state (storage).
                An expression contains constants, variables and function calls
                joined by operators. An expression could consist of a single
                variable in which case the value variable to which it is assigned
                can be thought of as a pointer to that variable (or rather the
                value of the variable, i.e. its output, since a variable cannot
                be assigned (changed) via a value variable pointing to it).

                The availability (see later) of a value variable depends on any
                queue reads in the expression it represents.

            \subsubsection{selectvalue variables}

                \index{variable!selectvalue}
                \index{mode!selectvalue}
                \index{selectvalue mode variables}
                Variables of mode {\bf selectvalue} represent buses or selectors.

                The availability (see later) of a selectvalue variable depends
                on any queue reads in the equivalent expression.

            \subsubsection{input variables}

                \index{variable!input}
                \index{mode!input}
                \index{input mode variables}
                Variables of mode {\bf input} are external inputs.

                An input variable is always (read) available.

            \subsubsection{output variables}

                \index{variable!output}
                \index{mode!output}
                \index{output mode variables}
                Variables of mode {\bf output} are external outputs.

                An output variable is always (write) available.

            \subsubsection{static variables}

                \index{variable!static}
                \index{static mode variables}
                Variables of mode {\bf static} have state - they are registers.
                They have no associated signaling. They may be global, or local to
                modules, procedures, functions or code blocks. Static
                variables involve no signaling and their state changes are not
                implicitly synchronised with pipeline operations except in as much
                as module code may exercise local control.

                Static variables containing arithmetic types may be any size
                but initialisation or assignment from immediate expressions or
                variables are restricted at the moment to a maximum of 64 bits for
                integer or 52.11 for floating point.

                A static variable is always available.

            \subsubsection{queue variables}

                \index{variable!queue}
                \index{mode!queue}
                \index{queue mode variables}
                Variables declared {\bf queue} are dataflow variables which are
                mainly used as pipeline connections between modules. They have
                associated signaling to indicate their availability for reading,
                examination or writing (assignment). When a datum is written to a
                queue, the datum becomes available at the output of the queue. A
                queue may contain more than one datum, in which case the data will
                be available in the order they were written. In a module the
                written end of a queue is referred to as an {\bf output queue}
                since it is assigned to. The read or examined end of a queue is
                referred to as an {\bf input queue}.

                An input queue is available when it has a datum which can be read
                or examined (a datum which is read is removed from the queue
                whereas a datum which is examined remains in place), i.e. it is
                not empty.

                An output queue is available when it can be written to, i.e. it is
                not full.

                Queue variables containing arithmetic types may be any size but
                initialisation or assignment from immediate expressions or
                variables is limited to a maximum of 64 bits for integer or
                52.11 for floating point.
                
                A queue is 'synchronous' if it is written and read in the same
                clock domain.
                
                A synchronous queue has a default buffer depth of 2. An
                asynchronous queue has a default buffer depth of 16. These
                depths can be explicitly set.
                
                \index{null!queue}
                \index{queue!null}
                A queue may be type "null". This is useful for signalling as
                the queue can be written (null must be assigned) or read
                but no actual data storage is used.

            \subsubsection{cmemory and rmemory variables}

                \index{variable!memory}
                \index{mode!memory}
                \index{memory mode variables}
                Variables of mode {\bf memory} represent RAM elements in the
                FPGA.

                Memory variables can only be accessed via inbuilt procedures and
                functions and hence availability does not apply as they cannot be
                used in expressions or assignment statements.

           \subsubsection{priority variables}\label{priorityvariables}

                \index{variable!priority}
                \index{mode!priority}
                \index{priority mode variables}
                Variables of mode {\bf priority} implement a priority encoder
                and are 
                used for arbitration. The type must be a one dimensional array
                of "log" (if the type is omitted "[20]log" will be used). The
                array members are assigned combinatorial values in the same way
                as value mode variables. However, the resultant values of the
                logical array are not simply the input expressions but are
                arbitrated values, the lowest subscript having highest priority.
                Where one or more input expressions to the priority array are
                true, one and only one array value will be true and that one
                will correspond to the true input with the lowest subscript.

                There is an inbuilt function, {\bf prielse()}, which returns
                the value of the else case, i.e. a value which is true when
                no priority inputs are true.

                The entire array need not be assigned. Unused entries will be
                ignored, but where an entry is evaluated (its output is used)
                it must have an assigned input value. The converse is not
                true - an entry that has been assigned need not be evaluated.
                Assignments need not be done in any particular order. This method
                of usage is designed to allow an array that is larger than
                required to be created after which entries can be
                assigned as immediate execution progresses. At the completion of
                immediate execution a priority encoder will be created in target
                code of a size which matches the entries in the array that have
                actually been used. The avoids prior knowledge of the required
                size of the encoder.
        
                A priority variable array member may be assigned more than once
                but the default behaviour is that the final value assigned prior
                to termination of immediate execution prevails, all other
                previous assignments being lost. This can be changed by setting
                attribute "multiplein" to true in which case all assignments
                will be collected and ORed together on completion of immediate
                execution to form the input to that priority array member.

                {\bf At the moment assignments to priority variables cannot
                contain queue reads. This restriction may be removed later.}

                A priority variable is always available.

            \subsubsection{clock variables}
            
                \index{variable!clock}
                \index{mode!clock}
                \index{clock mode variables}
                Variables of mode {\bf clock} are used to represent signals.
                There may be a number of clocks to which program logic is
                synchronised, but there will also be associated intermediate
                clock signals used to connect together the hardware components
                used to generate these program clocks, such as input buffers,
                delay-locked loops and clock buffers.


        \subsection{variable creation}

            \index{variable!creation}
            Variables are created during immediate execution by means of
            inbuilt procedures. The most general procedure used is {\bf
            var(...)}. This has two to four arguments -
            \begin{lstlisting}[gobble=12, language=threepl]
               var(name, mode, type, init)
            \end{lstlisting}

            {\bf name} is the identifier for the created variable.

            Most variable creation inbuilt procedures allow creation of multiple
            variables of the same type, mode and initial value by replacing the
            {\em name} argument by an identifier list enclosed in braces {\em \{v1,
            v2,\dots\}}.

            \index{variable!mode}
            The {\bf mode} argument is optional and must be one of the keywords
            {\bf immediate}, {\bf value}, {\bf selectvalue}, {\bf static} or
            {\bf queue} - memory or input mode variables must be created by
            other procedures (see below). If the mode argument is omitted the
            mode will be {\bf immediate}.

            \index{variable!type}
            The type argument must be a type-valued expression or a string-valued
            expression describing the type. A type-valued expression can be a
            variable or variable member of type "type", the inbuilt functions {\bf
            type()} or {\bf struct()}, or a user-defined function returning type
            "type".
            
            Type specifications in general and type variables in particular
            are described in a \autorefph{sec:types}.

            \index{variable!initialisation}
            \index{initialisation!variable}
            The final argument to {\bf var()} is optional and is an
            initialisation value. This may take one of three forms -
            \blist
            \item   a single immediate expression of appropriate type for a
                    variable of primitive type
            \item   a variable reference of matching type to the created
                    variable
            \item   a compound value
            \blend
            
            For the first two of the above cases where the mode is immediate the
            type argument may be the string "empty" in which case the variable type
            will be the type of the initialisation argument.
            
            Where initialisation is provided and more than one variable is being
            created (using \{...\}) then all variables will be initialised
            identically.

            Immediate, static and queue variables may be initialised. If they are a
            compound type the initialisation argument will usually be a compound
            value or a compound variable of suitable type, but a compound
            type may be initialised by a single immediate value in which case all
            members of the compound type will be initialised to that single value
            (the members of the compound type must all be the same primitive type
            appropriate to the initialisation value).

            Uninitialised immediate and static variables are automatically
            initialised to 0, 0.0, false or "" as appropriate to the type.

            Queue variables when initialised contain a single datum.
            Uninitialised queues have no value as they are not available.
            Initialised queues have a value and they are available.

            A value mode variable can be initialised, but this does not preclude
            assignment of one or more other values later. A value mode variable can
            be initialised with an immediate constant in which case it behaves as a
            target constant unless assigned another expression. Uninitialised value
            variables have no default value. If the value variable is a compound
            type it may be initialised by a compound value.

            Selectvalue variables cannot be initialised but their output when
            no input is selected can be specified.

            Examples of variable creation -
            \begin{lstlisting}[gobble=12, language=threepl]
               var(byte, "int", 8);
               var(sum, value, "uint:16");
               var(even, static, "log");
               var(pixel, queue, "uint:" + byte);
               var(trigger, queue, "null");
            \end{lstlisting}

            The first has no mode declared, so is immediate. It is given an
            initial value of 8. The remainder are target variables. The
            second last example uses the '+' operator to concatenate an
            integer to a string to form the word size. The final example is a
            queue with no data which can be used as an event trigger.
            These examples all use string-valued expressions for the
            type argument. The following examples illustrate the use of
            type-valued arguments -
            \begin{lstlisting}[gobble=12, language=threepl]
               var(t1, type, "int:8");
               var(v1, value, t1);
               var(s2, static, type(v1));
               var(t2, type, type(s2));
               var(q1, queue, "[4]t2");
               var(s3, static, "(f1, int:8, f2, [4]t1)");
            \end{lstlisting}

            \index{mode!immediate!creation}
            \index{variable!immediate!creation}
            For convenience the following procedures can be used instead of
            {\bf var(...)} to create primitive immediate variables -
            \begin{lstlisting}[gobble=12, language=threepl]
               int(name, init);
               log(name, init);
               str(name, init);
               file(name, init);
               type(name, typespec);
               struct(name, field1, type1, field2, type2, .....);
            \end{lstlisting}
            In each case a single variable of the type described by the
            procedure name is created. The initialisation arguments are optional.

            These are further desribed in 
            \blist
            \item   {\bf int()} (\autorefph{sec:ipint})
            \item   {\bf log()} (\autorefph{sec:iplog})
            \item   {\bf str()} (\autorefph{sec:ipstr})
            \item   {\bf file()} (\autorefph{sec:ipfile})
            \item   {\bf type()} (\autorefph{sec:iptype})
            \item   {\bf struct()} (\autorefph{sec:ipstruct})
            \blend

            \index{mode!target!creation}
            \index{variable!target!creation}
            Similarly target variables can be created by-
            \begin{lstlisting}[gobble=12, language=threepl]
                   static(name, type, init);
                   value(name, type, init);
                   queue(name, type, init);
            \end{lstlisting}
            which are equivalent to -
            \begin{lstlisting}[gobble=12, language=threepl]
               var(name, static, type, init);
               var(name, value, type, init);
               var(name, queue, type, init);
            \end{lstlisting}
            In all the above {\em init} is optional.

            These are further desribed in 
            \blist
            \item   {\bf static()} (\autorefph{sec:ipstatic})
            \item   {\bf value()} (\autorefph{sec:ipvalue})
            \item   {\bf queue()} (\autorefph{sec:ipqueue})
            \blend

            \index{mode!memory!creation}
            \index{variable!memory!creation}
            Memory variables are created by the inbuilt procedures -
            \begin{lstlisting}[gobble=12, language=threepl]
               cmemory(name, ports, atype, dtype, init);
               rmemory(name, ports, atype, dtype, init);
            \end{lstlisting}
            These create combinatorial output RAM (mode {\bf cmemory}) and
            registered output RAM (mode {\bf rmemory}) respectively. {\em name} is
            an identifier for the memory, {\em ports} is the number of ports, {\em
            dtype} is the data type, {\em atype} is the address type and {\em init}
            is an optional initialisation array (immediate array variable or
            compound value).
            
            See {\bf cmemory()} (\autorefph{sec:ipcmemory}) and
            {\bf rmemory()} (\autorefph{sec:iprmemory}),
                        
            \index{mode!clock!creation}
            \index{variable!clock!creation}
            Clock variables are created by the inbuilt procedure -
             \begin{lstlisting}[gobble=12, language=threepl]
                clock(name, ...attributes...);
            \end{lstlisting}
            
            See {\bf clock()} (\autorefph{sec:ipclock}).
           
            \index{mode!input!creation}
            \index{variable!input!creation}
            Input variables are created by the inbuilt procedure -
            \begin{lstlisting}[gobble=12, language=threepl]
               input(name, type, ...attributes...);
            \end{lstlisting}
            {\em name} is an identifier for the input variable.
            {\em type} is the variable type.
            {\em attributes} is optional and is either one or more key=value
            attribute arguments or a map of attributes.These arguments can be
            omitted and attributes added later by using inbuilt procedure {\bf
            attributes()}.
            
            See {\bf input()} (\autorefph{sec:ipinput}).
            
            \index{mode!output!creation}
            \index{variable!output!creation}
            Output variables are created by the inbuilt procedure -
            \begin{lstlisting}[gobble=12, language=threepl]
               output(name, type, source, attributes);
            \end{lstlisting}
            {\em name} is an optional identifier for the output variable.
            {\em type} is the variable type (also optional).
            {\bf source} is an optional source value for unconditional assignment.
            {\em attributes} is an optional map of attributes. The argument
            can be a map containing the attributes or can be a compound value
            containing the attributes as key=value entries in one-dimensional
            list. This argument can be omitted an attributes added later by
            using inbuilt procedure {\bf attributes()}.
            
            See {\bf output()} (\autorefph{sec:ipoutput}).
            
            \index{variable!scope}
            \index{scope!variable}
            Variables are normally created within the scope containing the call
            to the variable creation inbuilt procedure. However, if the variable
            identifier has a leading "global.", "file.", "calling." or "local."
            it is created in the specified scope.

            Using "global.", global variables can be created from within any
            scope.

            Using "file.", variables can be created in the current (deepest)
            file module scope.
           
            Using "local." allows the conditional creation of variables in
            the scope of a module, procedure or function. For example the
            code -
            \begin{lstlisting}[gobble=12, language=threepl]
               module
               m (...) {
                   if (doit)
                       static(s, "log");
               }
            \end{lstlisting}
            will create static variable {\bf s} in the module if variable
            {\bf doit} is true. However, the following code -
            \begin{lstlisting}[gobble=12, language=threepl]
               module
               m (...) {
                   if (doit) {
                       static(s, "log");
                       int(j);
                   }
               }
            \end{lstlisting}
            creates variables {\bf s} and {\bf j} in the scope of the
            {\bf if} block and not in the module scope. By changing
            the code to -
            \begin{lstlisting}[gobble=12, language=threepl]
               module
               m (...) {
                   if (doit) {
                       static(local.s, "log");
                       int(local.j);
                   }
               }
            \end{lstlisting}
            variables {\bf s} and {\bf j} are created in the scope of
            the called module.
            
            Using "calling." allows the creation of variables in
            the scope of the previous module, procedure or function.
            This is useful for writing procedures to create variables, in
            which case the created variable needs to be in the calling
            scope, not the scope of the called procedure.

            Variables may be created at any point in the code, but they
            must have been created prior to being referenced.

            A variable reference where the identifier has a scope prepended
            will attempt to locate the variable only in that scope.

            It is important to note that variables exist independently of the
            scope in which they are created. The scope only contains the map
            which ties the variable identifier to the variable, not the
            variable itself. That map disappears with the scope. If a pointer
            to a variable is kept then that variable can be referred to via
            the pointer even though the scope in which the variable was
            created has disappeared. For example the following code illustrates
            a variable created within a function which returns a pointer to that
            variable so that it can be used elsewhere -
            \begin{lstlisting}[gobble=12, language=threepl]
               var(p, "->");
               p = make_static("uint:12");

               when (l)
                   p-> <- v;

               function
               make_static (str(t)) {
                   static(s, t);
                   return(&s);
               }
            \end{lstlisting}
            Here a pointer variable {\bf p} is assigned a pointer to a
            static variable created by function {\bf make\_static}. It can
            then be used by dereferencing the pointer; in the example the
            static variable is conditionally assigned a value {\bf v} via
            the pointer {\bf p}.
            
            \index{variable!creation!computed identifier}
            \label{subsec:strvarcreat}
            When creating any variable the identifier may be computed
            by using the indirect operator on a string variable rather than
            using an explicit identifier, e.g.
            \begin{lstlisting}[gobble=12, language=threepl]
               int(i, 23);
               str(name, "global.var" + i);
               log(name->);
            \end{lstlisting}
            will create a logical immediate variable in global scope whose
            identifier is {\bf var23}.






    \section{Types}\label{sec:types}

        \index{type}
        \index{immediate!variable!type}
        Immediate variables may be of primitive type
        \blist
        \item   bits
        \item   uint
        \item   int
        \item   log
        \item   str
        \item   float
        \item   \verb'->'
        \item   type
        \item   list
        \item   map
        \item   file
        \item   empty
        \blend

        {\bf bits} is a long unsigned integer (64 bit). Arithmetic operators
        +, -, *, / and \verb'%' and relational operators \verb'<', \verb'<=',
        \verb'>' and  \verb'>=' cannot be used on this type, but the bitwise
        logical operators and equality and inequality operators can be used.

        {\bf int} is a signed long integer (64 bit).

        {\bf uint} is an unsigned long integer (64 bit).

        {\bf log} is a boolean type whose value is 'true' or 'false'.

        {\bf str} is a character string type.

        {\bf float} is a floating point double precision type.

        {\bf \verb'->'} is a pointer to a variable, or components of a variable.

        {\bf type} is a type specification.

        {\bf list} is a list.

        {\bf map} is a sorted map with string keys.

        {\bf file} is a file.

        {\bf empty} for immediate mode specifies that the type will be taken from
        the variable initialisation argument.

        Integer constants and expressions which are 0 or positive are automatically
        given type {\bf uint}. Integer variables can be declared {\bf uint} but
        they are implemented as {\bf int} and no checks are made that the value of
        a uint remains non-negative.

        \index{target!variable!type}
        Target variables may be of primitive type
        \blist
        \item    bits:{\em n}
        \item    int:{\em n}
        \item    uint:{\em n}
        \item    log
        \item    fixed:{\em m}.{\em n}
        \item    ufixed:{\em m}.{\em n}
        \item    float:{\em m}.{\em n}
        \item    null
        \blend

        {\bf bits:{\em n}} is a word of n bits. This is used for bit patterns that
        do not represent weighted binary arithmetic values, e.g. ASCII
        characters or gray codes. It can also be used where a variable is
        used at different times for different types, e.g. a communications
        bus in which case other types are cast in or out of it. This type
        cannot be used in arithmetic operations.

        {\bf int:{\em n}} and {\bf uint:{\em n}} are integers or unsigned
        integers respectively whose size in bits is explicitly
        declared by {\em n}. The size {\em n} must be an immediate expression.

        {\bf log} is a logical type whose value is {\bf true} or {\bf false}.

        {\bf fixed:{\em i}.{\em f}} and {\bf ufixed:{\em i}.{\em f}} are fixed
        point or unsigned fixed point types respectively with an {\em i} bit
        integer part and {\em f} bit fractional part. {\em i} or {\em f} may be
        zero (but not both). The maximum size of the fractional part is 63 bits.
        
        {\bf float:{\em e}.{\em f}} is a floating point type with an {\em f}
        bit mantissa and {\em e} bit biased exponent.

        \index{null!queue}
        \index{queue!null}
        A {\bf null} type can be used where a queue is being used only as an
        event trigger. It cannot be used for variables of other modes.

        \index{compound!type}
        \index{type!compound}
        Variables may be of compound type
        \blist
        \item   struct
        \item   array
        \blend

        A {\bf struct} is a compound type consisting of fields named by
        an identifier. The fields may be any type, primitive or compound.

        An array can be any type, including {\bf struct}s.
        
        There is an enumerated type. This is effectively a primitive type
        but unlike the others is not specified the same way (see below).

        \subsection{type specifications}
        
            \index{type!specification}
            Where variable creation procedures require a type argument
            this can either be a string type description expression or
            a type variable (\autorefph{sec:iptype}).
            
            If a string representation of a type is used, which is the most
            common way to specify types, this is just a string constant,
            a string expression or a string variable.
            
            A type is specified by one of the following -
            \blist
            \item   A type keyword with optional trailing ":" and width
                    where appropriate.
            \item   The identifier of a type variable.
            \item   The identifier of a string variable containing a type
                    description string.
            \item   A structure description.
            \blend

            Type keywords are
            \blist
            \item   {\bf bits} - a word of arbitrary bit encoding
            \item   {\bf uint} - an unsigned integer
            \item   {\bf int} - a signed integer
            \item   {\bf ufixed} - an unsigned fixed point value
            \item   {\bf fixed} - a fixed point value
            \item   {\bf float} - a floating point value
            \item   {\bf log} - a boolean
            \item   {\bf str} - a string
            \item   {\bf \verb'->'} - a pointer
            \item   {\bf type} - a type
            \item   {\bf class} - a class
            \item   {\bf list} - a list
            \item   {\bf map} - a map
            \item   {\bf file} - a file
            \item   {\bf null} - a dataless queue
            \item   {\bf empty} - the immediate type will be taken from the
                                  initialisation argument
            \blend
            An optional data width can only be appended to target mode numeric
            types ({\bf uint}, {\bf int}, {\bf fixed} or {\bf float)}). The
            width specification is separated from the type by a colon. For {\bf
            uint}, {\bf int} a single width in bits is required. The {\bf
            ufixed} and {\bf fixed} types have integer part width followed by
            '.' and then fractional part width. For type {\bf float} the width
            specification is given in the form {\bf {\em e}.{\em m}} where {\em
            m} gives the number of bits in the mantissa and {\em e} the number
            of bits in the biased exponent (the total size is 1+n+m as there is
            also a sign bit). The actual mantissa fraction has an implied
            integer 1 preceding the fraction.
            
            A type variable can be created by means of inbuilt procedures {\bf
            type()} or {\bf var()}. Inbuilt procedure {\bf struct()} also creates a
            type variable.
            
            Assuming we have created the variables -
            \begin{lstlisting}[gobble=12, language=threepl]
               int(i, 8);
               str(s, "uint:");
            \end{lstlisting}
            the following are valid type strings -
            \begin{lstlisting}[gobble=12, language=threepl]
               "uint:8"
               "uint:"+i
               s+8
            \end{lstlisting}
            The first is a string constant. The second and third are string
            expressions, the "+" giving concatenation when one operand is
            a string (the other being converted to a string if necessary).
            
            Any spaces in a type specification string are removed prior
            to its evaluation.
            
            \index{type!type}
            A type variable is a variable of type "type" and is created by
            the {\bf type()} inbuilt procedure, e.g.
            \begin{lstlisting}[gobble=12, language=threepl]
               type(t, "uint:"+8);
            \end{lstlisting}
            The {\bf type()} inbuilt procedure takes a type specification string
            and constructs a type variable containing that type value. There is an
            inbuilt function, {\bf typestring()}, which will return a type
            description string of the type of the argument. These roughly
            complementary facilities allow manipulation of types. The type of a
            module, procedure or function argument can be turned into a string,
            modified, then used as the type of a new variable. Similarly types can
            be constructed initially using parameters.

            The {\bf struct()} inbuilt procedure also creates a variable
            of type "type".
           
            The following example creates static variables a, b and c which
            are unsigned integers of width 8 bits -
            \begin{lstlisting}[gobble=12, language=threepl]
               int(i, 8);
               str(s, "uint:");
               var(a, static, "uint:8");
               var(b, static, "uint:"+i);
               var(c, static, s+8);
            \end{lstlisting}

            Note that types {\em uint} and {\em int} which can be used for
            both immediate and target mode variables {\bf must} have a width
            specified for target mode variables but {\bf must not} have a
            width specified for immediate variables.
            
            Where a type key word is expected an identifier of a variable of type
            "str" or type {\bf type} may be used. Such an embedded variable must
            not be followed by a width specification! If the embedded variable is
            type "str" then the string contained in the variable must itself be a
            valid type description string. The identifier may have a preceding
            scope such as {\bf calling.}, {\bf global.} etc.
            Thus type string expressions and type variables can be nested,
            e.g.
            \begin{lstlisting}[gobble=12, language=threepl]
               str(s1, "uint");
               int(i, 8);
               type(t, s1+i);
               queue(p, "[4][3]t");
               static(s, "[5]s1");
            \end{lstlisting}
            This nesting of string and type variables can be to any depth.
            The following is incorrect as an embedded variable cannot be followed
            by a width specification, i.e. it must be a complete type -
            \begin{lstlisting}[gobble=12, language=threepl]
               type(t, "uint");
               queue(p, "t:8"); // ERROR!
            \end{lstlisting}
            
            Note that variables of type {\bf type} can be assigned to variables of
            type "str" and vice versa and in each case a conversion will be
            applied. However a string assigned to a type variable must have valid
            type syntax or a fatal error occurs.
            
            Note also that a string describing a type which contains an embedded
            identifier must be able to access that identifier when the string
            is used to create a type. If for example the string is passed
            to a module, procedure or function and is used within that
            code body as the type of a variable then the variable embedded
            within the type string must be accessible in the current scope.
            Where a string is passed as an argument to a type parameter the
            conversion occurs in the scope of the call, but if the string
            is passed to a string parameter and then used later then there
            may be a problem with the scope.
            
            The following example works as expected -
            \begin{lstlisting}[gobble=12, language=threepl]
               struct(s,
                   f1, "int",
                   f2, "log"
               );
               p("[3]s");

               procedure p(type(t)) {
                   var(v, t);
                   .
               }
            \end{lstlisting}
            The argument is a string describing a type and contains the embedded
            type variable {\bf s} (this could also have been a string variable).
            The procedure parameter {\bf t} is type {\bf type} and will be
            assigned the type derived from the string. When the argument is
            assigned to the procedure parameter the embedded
            variable {\bf s} will be evaluated in the calling scope.
            
            If however we change the example to -
            \begin{lstlisting}[gobble=12, language=threepl]
               struct(s,
                   f1, "int",
                   f2, "log"
               );
               p("[3]s");

               procedure p(str(t)) {
                   var(v, t);  // cannot find embedded variable 's'!
                   .
               }
            \end{lstlisting}
            then the argument is passed unchanged and an attempt to use it as
            a type within the procedure body fails as embedded variable {\bf s}
            is not accessible within the scope of the procedure body.
            
            There are two ways this error could be overcome. We could
            explicitly specify the calling scope in the argument -
            \begin{lstlisting}[gobble=12, language=threepl]
               struct(s,
                   f1, "int",
                   f2, "log"
               );
               p("[3]calling.s");

               procedure p(str(t)) {
                   var(v, t);
                   .
               }
            \end{lstlisting}
            but this seems rather unsatisfactory. Alternatively we could force the
            conversion from string to type by using inbuilt function {\bf
            strtotype()} \autorefph{sec:ifstot}. The parameter is of type "str" so
            the assignment of the resulting type argument to the string parameter
            would convert the value back to string automatically.
            \begin{lstlisting}[gobble=12, language=threepl]
               struct(s,
                   f1, "int",
                   f2, "log"
               );
               p(strtotype("[3]s"));

               procedure p(str(t)) {
                   var(v, t);
                   .
               }
            \end{lstlisting}
            We could have also used an intermediate type variable -
            \begin{lstlisting}[gobble=12, language=threepl]
               struct(s,
                   f1, "int",
                   f2, "log"
               );
               type(it, "[3]s");
               p(it);

               procedure p(str(t)) {
                   var(v, t);
                   .
               }
            \end{lstlisting}
            


        \subsection{arrays}

            \index{array}
            \index{type!array}
            \subsubsection{creating arrays}   

                \index{type!array!creation}
                \index{array!creation}
                Arrays are created by prepending dimensions in brackets
                to the type, e.g.
                \begin{lstlisting}[gobble=16, language=threepl]
                   var(pix, value, "[3][3]uint:8");
                   var(sum, static, "[10]int:16");
                   var(q1, queue, "[2]uint:8");
                \end{lstlisting}
                \index{array!dimensions}
                \index{array!undimensioned}
                \index{dimensions!missing}
                One or more dimensions may be omitted (empty brackets) for module,
                procedure or function parameters in which case they will be
                supplied from the matching argument. A parameter array which is
                initialised may also have some dimensions omitted, in which case
                the missing dimensions will be inferred from the  argument if
                supplied or else by default from the initialisation. Either or both
                the argument or the initialisation may be a variable of matching
                compound type or a compound value.
                
                If a parameter has one or more undimensioned arrays and it has
                no initialiser and is not passed an argument, the undimensioned
                arrays will be given dimensions of 0.

                Array dimensions cannot be omitted where a variable is not a
                module, procedure or function parameter and is not initialised.
                
                The array type may be an identifier of a type or string variable
                which specifies any type including an array.

                Arrays may be used for modes {\bf immediate}, {\bf value}, {\bf
                selectvalue}, {\bf static} or {\bf queue}. Variables of modes
                {\bf memory} or {\bf clock} do not have types and hence cannot
                be created as array types. Variables of mode {\bf priority}
                must be arrays of {\em log} (default type if omitted).

                Where array dimensions are given by immediate variables or
                expressions a type string is created by concatenating the
                dimensions and surrounding brackets, e.g.
                \begin{lstlisting}[gobble=16, language=threepl]
                   int(dim1, 8);
                   int(dim2, 4);
                   var(a, queue, "["+dim1+"]["+dim2+"]uint:8");
                \end{lstlisting}
                Care should be taken to parenthesize nested arithmetic
                expressions where necessary to ensure that they are evaluated
                first into numeric values prior to concatenation with strings,
                otherwise the evaluation order may be erroneous, for example
                \begin{lstlisting}[gobble=16, language=threepl]
                   int(i, 2);
                   int(j, 3);
                   var(a, queue, "["+i+j+"]uint:8");
                \end{lstlisting}
                will give a dimension of 23 not 5.

            \subsubsection{referencing array members}

                \index{type!array!reference}
                \index{array!reference}
                \index{array!subscripts}
                An array subscript may be -
                \blist
                \item    a single subscript
                \item    a subscript range
                \item    missing
                \blend
                Trailing subscripts may be omitted entirely, implying that the
                array reference includes the full dimensions where the
                subscripts are omitted.
                
                A subscript range selects a sub-array.
                If the first subscript of the pair is higher than the second
                the sub-array members will be in reverse order.

                Where arrays have greater than one dimension subscripts are
                concatenated, each subscript being enclosed in brackets.
                Intermediate subscripts may be omitted by using empty brackets.
                Subscripts or subscript ranges may be any expression whose type
                is {\em uint} or {\em int}. Subscripts may be immediate, value,
                selectvalue, static or queue mode.
                             
                The following are subscript examples -
                \begin{lstlisting}[gobble=16, language=threepl]
                   a[2]        constant subscript
                   a[i]        variable subscript
                   a[2..j-1]   subscript range
                   a[]         missing subscript - equivalent to full range
                   a[2][i]
                   a[2..4][i..i+3]
                \end{lstlisting}

                Some caution should be exercised when using target subscripts.
                These employ considerable logic but more importantly the logic
                delays may prove difficult to fit within timing constraints. In
                particular delays can become large where the array dimensions
                are large or array references are nested. At low clock rates
                this is not a concern.



        \subsection{structs}\label{sec:ssecstructs}
        
            \index{struct}
            \index{type!struct}
            Structures are usually created by using inbuilt procedure
            {\bf struct()}, but they may also be specified by
            a type string.

            \subsubsection{using the struct() procedure}
            
                \index{type!struct!creation}
                \index{struct!creation}
                The first argument of the inbuilt procedure {\bf struct(...)}
                is an identifier giving the struct name which will be created
                as a variable of type "type". This is followed by a number of
                argument pairs, the first component being an identifier giving
                the field name and the second an associated type string-valued
                expression.
                
                If the first component of a pair is the keyword
                {\bf extends} then the second component must be a struct name.
                The fields of that struct will be included in this struct at
                that point. The type of a field can of course be another struct
                in which case that struct is nested (rather than being
                extended).
                
                The first component of a pair may be a comma-separated list
                of identifiers enclosed in braces in which case several fields
                of the same type will be created, one for each field identifier
                in the list.
                
                The second component of a pair is the type of the field. This
                is a string specifying the type or is a type variable.
                
                A field type cannot be, or contain, an undimensioned array.
                
                Fields are stored in the order in which they are listed.
                
                Examples -
                \begin{lstlisting}[gobble=16, language=threepl]
                   struct(pixel,
                       pix,    "uint:8",
                       hr,     "log",
                       vr,     "log"
                   );
                \end{lstlisting}
                This declares a pixel structure which has an 8-bit unsigned
                integer field called "pix', and horizontal and vertical reset
                indicators which are logical (1 bit). The arguments are all
                string constants but could be string variables or expressions.
                \begin{lstlisting}[gobble=16, language=threepl]
                   struct(win,
                       pixwin, "[3][3]pixel",
                       edge,   "log"
                   );
                \end{lstlisting}
                This example contains a field which has as its type the
                previous structure. If we create a variable of this type -
                \begin{lstlisting}[gobble=16, language=threepl]
                   var(a, queue, win);
                \end{lstlisting}
                then we can access the pixel value of one member of the window
                by -
                \begin{lstlisting}[gobble=16, language=threepl]
                   ...a[1][0].pix...
                \end{lstlisting}

                In the following case
                \begin{lstlisting}[gobble=16, language=threepl]
                   struct(pixaddr,
                       extends,  pixel,
                       address,  "uint:19"
                   };
                \end{lstlisting}
                we have created a structure which includes the fields of struct
                {\em pixel}. This is a single level struct and is not nested.

                To create multiple fields of the same type a list of field
                identifiers may be used, enclosed in braces, e.g.
                \begin{lstlisting}[gobble=16, language=threepl]
                   struct(pixaddr,
                       f1,          "int",
                       {f2, f3},    "log"
                   };
                \end{lstlisting}

                Structs are created within the scope containing the call to the
                structure creation inbuilt procedure unless the struct identifier
                has a leading "global.", "file.", "calling." or "local." in
                which case it is created in the specified scope. Structs may be
                created at any point in the code, but they must have been
                created prior to being referenced.

                A struct reference where the identifier has a leading "global.",
                "file.", "calling." or "local." will attempt to locate the
                struct only in the specified scope.

                Note that structs must not mix types {\bf uint}, {\bf int}, {\bf
                ufixed}, {\bf fixed} or {\bf float} such that some have a width
                specified and some do not. If the struct is to be used for
                immediate variables then there must be no widths specified. If the
                struct is to be used for target variables then widths must be
                given for these types. Types {\bf log} and {\bf enum} can be
                interpretted as either immediate or target mode so can be included
                in an immediate struct or a target struct. A target type can
                howvere be included in an immediate struct by using a pointer to
                the target type.

            \subsubsection{using a string specification}
            
                A struct can also be created using a string specification. This
                consists of outside parentheses enclosing field/type pairs, for
                example
                \begin{lstlisting}[gobble=16, language=threepl]
                   "(a, log, b, [3]uint:8)"
                \end{lstlisting}
                Field identifiers cannot be enclosed in braces; if several
                successive fields have the same type they must still be
                specified separately. String or type variable identifiers can
                be included where a type would be expected. The following two
                statements are equivalent -
                \begin{lstlisting}[gobble=16, language=threepl]
                   type(t, "(a, log, b, log, c, [3]uint:8)");
                \end{lstlisting}
                \begin{lstlisting}[gobble=16, language=threepl]
                   struct(t,
                       {a, b}, "log",
                       c       "[3]uint:8"
                   );
                \end{lstlisting}

                Any spaces in the specification string are removed prior
                to its evaluation.

                Structs may be used for modes {\bf immediate}, {\bf value},
                {\bf selectvalue}, {\bf static} or {\bf queue}. Structs cannot
                be used for modes {\bf memory}, {\bf clock} or {\bf priority}.

            \subsubsection{referencing struct members}
                
                \index{type!struct!reference}
                Struct members are accessed by using the struct name followed
                by "." and the field identifier, e.g.
                \begin{lstlisting}[gobble=16, language=threepl]
                   v.f1
                \end{lstlisting}
                If the field has a compound type the field name may be
                followed by more field identifiers or subscripts, e.g.
                \begin{lstlisting}[gobble=16, language=threepl]
                   v.f1[2].f2
                \end{lstlisting}
                
                Note that structs are interchangeable if their fields and
                associated types are identical. Thus the following code
                is legitimate -
                \begin{lstlisting}[gobble=16, language=threepl]
                   struct(s1, f1, "int", f2, "log"};
                   struct(s2, f1, "int", f2, "log"};
                   var(v1, s1);
                   var(v2, s2);
                   v1 = v2;
                \end{lstlisting}
                though there seems no reason to do this and the practice
                is confusing and to be discouraged.


        \subsection{lists}\label{sec:sseclists}
        
            \index{list}
            \index{type!list}
            Lists are created by inbuilt procedure {\bf var()} with type string
            "list". An initialiser may be provided which must be another list
            variable or a compound value, in which case values will be copied
            into the new list. If multiple list variables are created using a
            variable list as the first argument then independent lists will be
            created (they are not pointers to the same list). If an initialiser
            is used when creating multiple list variables, each will be
            initialised with data copied from the initialising list. If the
            initialiser is a compound value the list must not contain nested
            lists, null entries or key=value entries.
            
            List entries are always values, they are not pointers to other
            variables!     
            
            List entries must only be immediate, clock, cmemory or rmemory mode.
            Variables of other modes cannot be assigned to a list but pointers to
            them can be. \footnote {Those target modes are allowed as they do
            not have types and hence do not have subscripts or fields. The
            complexities of nested subscripts and fields preceding and following
            a list variable with its subscript (index) have been managed in the
            case of immediate mode variables but this has yet to be expanded
            to target modes. A future version of 3PL may allow all variable
            modes to be appended to lists.}
            
            List entries may be any type and entries do not have to be
            the same type.
            
            List operations are summarised in the following table -
            
            \begin{tabular}{| l | l |}
            \hline
            v = get(l, index);          & get the list entry \\
            .... l[index]....           & get the list entry \\
            append(l, index, v);        & append an item at the index position \\
            appendall(l, index, l2);    & append contents of another list \\
            change(l, index, v);        & change a list entry \\
            l[index] = v;               & change a list entry \\
            remove(l, index);           & remove a list entry \\
            n = entries(l)              & get the number of list entries \\
            \hline
            l = l2;                     & assign another list \\
            l = \{\dots\};              & copy a compound list \\
            l = null;                   & clear the list \\
            \hline
            \end{tabular}

            The inbuilt function {\bf entries()} returns the number of elements in
            a list (or a map).
            
            Lists may be nested within arrays, structs and other lists and maps.
            
            Lists may contain arrays, structs and other lists and maps.
            
            Lists may contain null entries.
            
            A list can referenced via a pointer and a list entry can be selected
            using a subscript after the pointer dereference.
            
            {\bf v = l[index] - get a list element using subscript notation}\\
            {\bf v = get(l, index) - get a list element using a function}
            
            The form using subscript notation was a later addition to 3PL.
            The inbuilt function has been retained as an alternative.
            
            List items are indexed starting at 0. The index may range from 0 to
            {\bf entries(l) - 1} inclusive. If {\bf index} is outside this range
            a fatal error occurs.
            
            A list item can also be accessed using a subscript. This can be done
            to any depth. If a list entry contains structs, arrays, lists or maps
            nested to any depth then trailing subscripts and fields can be used.
            In the case of lists the subscript again is the index into the list.
            In the case of maps the subscript may be a string (the key) or an
            integer (an index to a key).
            
            In the case of the following -
            
            \begin{lstlisting}[gobble=12, language=threepl]
               struct (s,
                   f1, "int",
                   f2, "[4]list"
               );
               var(v, s);
               var(va, "[3]int", {4, 6, 8});
               append(v.f2[1], 0, va);
            \end{lstlisting}
            we could access the nested integer array by -
            \begin{lstlisting}[gobble=12, language=threepl]
               ........v.f2[1][2]....
            \end{lstlisting}
            If we wanted the whole array then -
            \begin{lstlisting}[gobble=12, language=threepl]
               ........v.f2[1]....
            \end{lstlisting}
            would be appropriate.
            
            Note that where lists are assigned nested types the values are copies,
            not pointers to the original variable.
           



            {\bf append(l, index, v) - append a list element}
            
            The item {\bf v} is appended at the position indicated by {\bf
            index} in list {\bf l}. The index may range from 0 to {\bf
            entries(l)} inclusive. If {\bf index} is outside this range a fatal
            error occurs. An index of 0 inserts at the beginning of the list and
            an index equal to the list size appends to the end of the list. If
            the index argument is omitted the item will be appended to the end
            of the list. If the last argument is omitted null will be appended.


            {\bf appendall(l, index, l2) - append list elements from another list}
            
            The contents of list {\bf l2} are appended at the position indicated
            by {\bf index} in list {\bf l}. The index may range from 0 to {\bf
            entries(l)} inclusive. If {\bf index} is outside this range a fatal
            error occurs. An index of 0 inserts at the beginning of the list and
            an index equal to the list size appends to the end of the list. If
            the index argument is omitted the items will be appended to the end
            of the list.


            {\bf l[index] = v - change a list element using subscript notation}\\
            {\bf change(l, index, v) - change a list element using an inbuilt procedure}
            
            The form using subscript notation was a later addition to 3PL.
            The inbuilt procedure has been retained as an alternative.

            The item in list {\bf l} at position {\bf index} is replaced by item
            {\bf v}. The index may range from 0 to {\bf entries(l) - 1}
            inclusive. If {\bf index} is outside this range a fatal error
            occurs. The new value does not have to be the same type as the
            previous value. If the last argument is omitted the list item will
            be changed to null.
            
            A list item may be changed by assigning using a subscript, the
            subscript being the index into the list of the item to be changed.
            Where an item in a list has nested structs, arrays, lists or maps
            trailing subscripts and fields may be used to change items within
            those nested types.
            
            In the case of the following -
            
            \begin{lstlisting}[gobble=12, language=threepl]
               struct (s,
                   f1, "int",
                   f2, "[4]list"
               );
               var(v, s);
               var(va, "[8]int");
            \end{lstlisting}
            we could append an integer 3 to the nested list in {\em v} by -
            \begin{lstlisting}[gobble=12, language=threepl]
               append(v.f2[1], 0, 3);
            \end{lstlisting}
            We could then change that item to the value of variable {\em va}
            by -
            \begin{lstlisting}[gobble=12, language=threepl]
               v.f2[1] = va;
            \end{lstlisting}
            We could now assign to the nested integer array by -
            \begin{lstlisting}[gobble=12, language=threepl]
               v.f2[1][3] = 15;
            \end{lstlisting}


            {\bf remove(l, index) - remove a list element}

            The item in list {\bf l} at position {\bf index} is removed.
            The index may range from 0 to {\bf entries(l) - 1}
            inclusive. If {\bf index} is outside this range a fatal error
            occurs.
            


            {\bf l = l2 - assign a list}

            If the LHS of an assignment is a list variable and the RHS is a
            list, the LHS list will be replaced by the RHS list. Unlike other
            immediate assignments the LHS is assigned the RHS value, not a copy
            of the RHS value. If a copy is required then the inbuilt function
            {\bf copy()} \autorefph{sec:ifcopy} can be used.

            The RHS may be a one-dimensional compound value list, in which case it
            must not contain nested lists, null entries or key=value entries.
            The LHS list is cleared and values from the compound list are
            copied into the LHS list.


            {\bf l = null - clear a list}

            If the LHS of an assignment is a list variable and the RHS is
            explicitly the keyword {\bf null} (not a null value) the list is
            cleared.






        \subsection{maps}\label{sec:ssecmaps}
        
            \index{map}
            \index{type!map}
            Maps are created by inbuilt procedure {\bf var()} with type string
            "map". An initialiser may be provided which must be another map
            variable, or a compound value, in which case values will be copied
            into the new map. If multiple map variables are created using a
            variable list as the first argument then independent maps will be
            created (they are not pointers to the same map). If an initialiser
            is used when creating multiple map variables, each will be
            initialised with data copied from the initialising map. If the
            initialiser is a compound value the list must not contain nested
            lists or null entries and all entries must have key=value format.
            
            Map keys are strings (type "str"). The elements in a map are sorted
            in the lexicographic order of their keys and remain so following
            additions or deletions.
            
            Map entries must only be immediate, clock, cmemory or rmemory mode.
            Variables of other modes cannot be assigned to a map but pointers to
            them can be. \footnote {Those target modes are allowed as they do
            not have types and hence do not have subscripts or fields. The
            complexities of nested subscripts and fields preceding and following
            a map variable with its subscript (key) have been managed in the
            case of immediate mode variables but this has yet to be expanded
            to target modes. A future version of 3PL may allow all variable
            modes to be included in maps.}
            
            Map entries may be any type and entries do not have to be
            the same type.
            
            Map operations are summarised in the following table -

            \begin{tabular}{| l | l |}
            \hline
            v = get(m, key);    & get the map entry from a key \\
            v = m[key];         & get the map entry from a key \\
            v = get(m, index);  & get the map key at index \\
            v = m[index];       & get the map key at index \\
            put(m, key, v);     & put an item with a key \\
            m[key] = v;         & put an item with a key \\
            putall(m, m2);      & put items from another map \\
            remove(m, key);     & remove a map entry \\
            n = entries(m)      & get the number of map entries \\
            \hline
            m = m2;             & assign a map \\
            m = \{\dots\};      & copy a compound list \\
            m = null;           & clear the map \\
            \hline
            \end{tabular}

            The inbuilt function {\bf entries()} returns the number of elements in
            a map (or a list).

            
            {\bf v = m[key] - get a map element using a key (subscript notation)}\\
            {\bf v = get(m, key) - get a map element using a key (inbuilt function)}\\
           
            The map item for the key is returned.
            If no item with that key exists then null is returned.
            
            The use of a key as a subscript to access a map was a later addition
            to 3PL. The inbuilt function has been retained as an alternative.
            
            A map item may also be referenced by using the string key
            as a 'subscript'. If the map is nested within an array, struct, list
            or map additional subscripts, fields, indices or keys as appropriate may
            precede the map key itself. If the map item is itself a nested
            type containing arrays, structs, lists or maps then trailing
            subscripts, fields, indices or keys as appropriate can be appended.

            
            {\bf v = m[index] - get a map key using an index (subscript notation)}\\
            {\bf v = get(m, index) - get a map key using an index (inbuilt function)}
            
            The use of an integer index as a subscript to access a map was a later addition
            to 3PL. The inbuilt function has been retained as an alternative.
            
            Map items are indexed starting at 0. Map entries are ordered
            according to the lexicographic order of the keys. The index may
            range from 0 to {\bf entries(l) - 1} inclusive. If {\bf index} is
            outside this range a fatal error occurs. The key at that position is
            returned. To get a value instead of a key using an index the calls
            must be nested, e.g. get(get(m, index)) or alternatively m[m[index]].
            
            Access by means of an index allows iterative traversal of a whole map,
            e.g.
            \begin{lstlisting}[gobble=12, language=threepl]
               var(m, "map");
               .
               .
               .
               for(int(i) ; i<entries(m) ; i++)
                   println(m[m[i]]);
            \end{lstlisting}
            
            The items will be accessed in lexicographic order.


            {\bf m[key] = v) - put an element into a map using subscript notation}\\
            {\bf put(m, key, v) - put an element into a map using an inbuilt procedure}
            
            The item {\bf v} is placed in the map {\bf m} against key
            {\bf key}. If that key is already in the map the value will
            be replaced. If the last argument is omitted null will be put
            against the key. The key should be type "str" but if it is type
            "int" or type "uint" it will be converted to a string when
            the entry is inserted into the map.
            
            When using subscript form, subscripts, fields, indices and keys
            may precede or follow the key 'subscript' as appropriate where
            the map is itself is nested within a struct, array, list or map or
            the map value is itself of nested types.


            {\bf putall(m, m2) - put all elements from another map into a map}
            
            All items from map {\bf m2} are added to map {\bf m}. Where a
            key already exists its value will be replaced.

            
            {\bf v = remove(m, key) - remove a map element using a key}
            
            The map item for the key is removed from the map.
            If no item with that key exists then no change occurs.
            
            
            {\bf m = m2 - assign a map}
            
            If the LHS of an assignment is a map variable and the RHS is another
            map, the RHS map is assigned to the LHS variable. Unlike other
            immediate assignments the LHS is assigned the RHS object, not a copy
            of the RHS object. Thus m and m2 contain the same map and changes to
            map entries  via either variable will be reflected in the other. If
            a copy is required then the inbuilt function {\bf
            copy()} \autorefph{sec:ifcopy} can be used.      

            If the LHS of an assignment is a map variable and the RHS is a
            compound value, the LHS map is cleared and the values in the RHS
            are entered. The compound value list must not contain nested lists
            or null entries and all entries must be key=value format.


            {\bf m = null - clear a map}

            If the LHS of an assignment is a map variable and the RHS is
            explicitly the keyword {\bf null} (not a null value) the map is
            cleared, i.e. all entries are removed.
            
            Maps may be compared using the \verb'==' operator. Two maps
            will appear equal if they are indeed the same map (one assigned to
            the other) or if the maps are independent but represent the
            same mappings.



        \subsection{enums}\label{sec:ssecenums}
        
            An enumerated type is usually created by using inbuilt procedure
            {\bf enum()} or {\bf enumv()}, but it may also be specified
            by a type string. Enumerated types have ordered values which are
            selected by an associated identifier. Each value has an
            associated unique ordinal which is not related to the ordering.
            The order of enumerated type values is that of their occurrence
            in the argument list of the procedures which create the type.

            \index{type!enum!creation}
            \subsubsection{using the enum() and enumv() procedures}
                
                The first argument of the inbuilt procedure {\bf enumv(...)}
                is an identifier giving the enum name which will be created
                as a variable of type "type". This is followed by a number of
                argument pairs, the first component being an identifier
                and the second an associated ordinal value which must be
                a positive integer expression.  

                Example -
                \begin{lstlisting}[gobble=16, language=threepl]
                   enumv(e1,
                       none,   3,
                       typea,  7,
                       typeb,  5
                   );
                \end{lstlisting}
                The identifiers and ordinals must be unique within this
                enum, but the identifiers may be used in other enums
                or as variable identifiers without conflict.
                
                If specific ordinals are not required the simpler procedure
                {\bf enum()} can be used, e.g.
                \begin{lstlisting}[gobble=16, language=threepl]
                   enum(e1, none, typea, typeb);
                \end{lstlisting}
                which is the same as the previous example except that the
                ordinals are now 0, 1, 2 etc.
                
                
                
            \subsubsection{using a string specification}
            
                An enum type can also be created using a string specification.
                This consists of outside braces enclosing identifier/ordinal
                pairs, for example
                \begin{lstlisting}[gobble=16, language=threepl]
                   "{none, 3, typea, 7, typeb, 5}"
                \end{lstlisting}
                The ordinals must be explicitly specified.
                
                Note that
                \begin{lstlisting}[gobble=16, language=threepl]
                   type(e1, "{none, 3, typea, 7, typeb, 5}");
                \end{lstlisting}
                is equivalent to
                \begin{lstlisting}[gobble=16, language=threepl]
                   enumv(e1, none, 3, typea, 7, typeb, 5);
                \end{lstlisting}

                Any spaces in the specification string are removed prior
                to its evaluation.


            \subsubsection{selecting enum values}
            
                \index{type!enum!reference}
                An enum value is selected by appending a "." and
                the value identifier to the type identifier, e.g.
                \begin{lstlisting}[gobble=16, language=threepl]
                   type(e1, "{none, 3, typea, 7, typeb, 5}");
                   var(v, e1);
                   if (v == e1.typea)
                       v = e1.typeb;
                \end{lstlisting}
                
                Note that two enum types are interchangeable if their
                identifiers and ordinals are identical ordered lists.
                Thus for example if we create identical types -
                \begin{lstlisting}[gobble=16, language=threepl]
                   type(e1, "{none, 3, typea, 7, typeb, 5}");
                   type(e2, "{none, 3, typea, 7, typeb, 5}");
                \end{lstlisting}
                then types e1 and e2 are interchangeable, so we could write
                \begin{lstlisting}[gobble=16, language=threepl]
                   var(v, e1);
                       v = e2.typea
                \end{lstlisting}
                without error. However, there seems no legitimate reason to do
                this and to avoid confusion it is to be discouraged.
                
                
                
            \subsubsection{associated functions}
            
                \index{type!enum!functions}
                There are a number of inbuilt functions associated
                with enumerated types -

                \begin{tabular}{| l | l |}
                \hline
                ord(v)     \autorefph{sec:iford}     & the ordinal value of variable v \\
                \hline
                next(v)    \autorefph{sec:ifnext}    & the next value of variable v \\
                \hline
                prev(v)    \autorefph{sec:ifprev}    & the previous value of variable v \\
                \hline
                hasnext(v) \autorefph{sec:ifhasnext} & true if variable v has a next
                                             value (i.e. it is not the last
                                             value) \\
                \hline
                hasprev(v) \autorefph{sec:ifhasprev} & true if variable v has a previous
                                             value (i.e. it is not the first
                                             value) \\
                \hline
                first(x)   \autorefph{sec:iffirst}   & the first ordered value of
                                             variable or enum type x \\
                \hline
                last(x)    \autorefph{sec:iflast}    & the last ordered value of
                                             variable or enum type x \\
                \hline
                \end{tabular}

                Note that the order of enum values is not given by the
                magnitude of the ordinals but by their order in the argument
                list of the creation procedure.
                
                The identifier of an enum value can be obtained by using
                the inbuilt function {\bf tostring()} \autorefph{sec:iftostring}.
  

        \subsection{files}

            \index{files}
            Procedure {\bf println()}(\autorefph{sec:ipprintln}) can be used
            to write to standard output.

            A file variable is created using the inbuilt procedure {\bf file()}.
            To read or write a file it must be opened. This is done by the
            inbuilt procedure {\bf open()}(\autorefph{sec:ipopen}), e.g.
            \begin{lstlisting}[gobble=12, language=threepl]
               file(f1, "file1");
               open(f1, "r");
            \end{lstlisting}
            where the "r" specifies opened for reading.           
                    
            A number of files are already open as global variables. {\bf stdin} is the
            standard input, {\bf stdout} and {\bf stderr} are standard output and standard
            error, and {\bf rpt} is the 3PL report file. The report file is {\bf
            design\_name.rpt} if 3PL was called with the -r option, otherwise it is {\bf
            stdout}.
            
            A file variable can be copied to another file variable, e.g.
            \begin{lstlisting}[gobble=12, language=threepl]
                file(f1, "file1");
                file(f2, "file2");
                f2 = f1;
            \end{lstlisting}
            after which f1 and f2 now share the same file descriptor.
            
            A file opened for reading can be read by inbuilt function {\bf
            readln()}(\autorefph{sec:ifreadln}). This function has one argument,
            which is the file variable, and returns a line from the file as a
            string.
            
            Formatted input lines can be read by
            {\bf scanf()}(\autorefph{sec:ipscanf})
            or {\bf fscanf()}(\autorefph{sec:ipfscanf}).
            
            If a file read is attempted at the end-of-file position an empty string
            is returned. This can be tested by inbuilt function
            {\bf hadeof()}(\autorefph{sec:ifhadeof}) which returns true if the
            previous read occurred at end-of-file.
            
            A line can be written to a file opened for writing by
            inbuilt procedure {\bf writeln()}(\autorefph{sec:ipwriteln}) where the
            first argument is the file variable and the second argument is the
            string to be written.
            
            Formatted lines can be written using {\bf printf()}(\autorefph{sec:ipprintf})
            or {\bf fprintf()}(\autorefph{sec:ipfprintf}).
            
            On completion the file can be closed by {\bf
            close()}(\autorefph{sec:ipclose}) where its argument is the file
            variable.
            
            There are a number of inbuilt functions which can be used to test
            files and directories. These do not require file variables but
            operate on path names.

            \begin{tabular}{| l | l |}
            \hline
            {\bf filecanread()}(\autorefph{sec:iffilecanread})            & file can be read\\
            \hline
            {\bf filecanwrite()}(\autorefph{sec:iffilecanwrite})          & file can be written\\
            \hline
            {\bf filedelete()}(\autorefph{sec:iffiledelete})              & delete file\\
            \hline
            {\bf fileexists()}(\autorefph{sec:iffileexists})              & file exists\\
            \hline
            {\bf filegetparent()}(\autorefph{sec:iffilegetparent})        & absolute path above file\\
            \hline
            {\bf fileisdir()}(\autorefph{sec:iffileisdir})                & file is a directory\\
            \hline
            {\bf fileisfile()}(\autorefph{sec:iffileisfile})              & file is a file\\
            \hline
            {\bf filelastmod()}(\autorefph{sec:iffilelastmod})            & time of last modification\\
            \hline
            {\bf filelength()}(\autorefph{sec:iffilelength})              & file length\\
            \hline
            {\bf filelist()}(\autorefph{sec:iffilelist})                  & list of files and sirectories\\
            \hline
            {\bf filelistdirs()}(\autorefph{sec:iffilelistdirs})          & list of directories\\
            \hline
            {\bf filelistfiles()}(\autorefph{sec:iffilelistfiles})        & list of files\\
            \hline
            {\bf filemkdir()}(\autorefph{sec:iffilemkdir})                & create a directory\\
            \hline
            {\bf filerename()}(\autorefph{sec:iffilerename})              & rename a file or directory\\
            \hline
            {\bf filesetreadonly()}(\autorefph{sec:iffilesetreadonly})    & change permission to read-only\\
            \hline
            \end{tabular}
     
     



        \subsection{type conversion}

            \index{type!conversion}
            Types may be converted using the {\bf cast()} inbuilt function.

            \subsubsection{immediate type conversion}

                For immediate types, conversions between types {\em bits}, {\em
                uint}, {\em int} and {\em str} are allowed with the provisos -
                \blist
                \item   int to uint gives an error if the integer is -ve
                \item   str to numeric gives an error if not correctly formatted
                \blend
                A string to be converted to numeric must be
                \blist
                \item   a string of 0's and 1's preceded by "0b" or 0B"
                \item   a string of hexadecimal digits preceded by "0x" or "0X"
                \item   a string of octal digits preceded by "0"
                \item   a string of decimal digits whose 1st digit is not 0
                \blend
                An {\em int} may be converted to a {\em float} and vice versa.

            \subsubsection{target type conversion}

                Target mode type conversions allowed are shown in the following
                table -

                \begin{tabular}{| l | l || l | l | l | l | l | l | l | l |}
                \hline
                \multicolumn{2}{| c |}{} & \multicolumn{8}{| c |}{to} \\
                \multicolumn{2}{| c |}{} & bits & uint & int & ufixed & fixed & log & float & enum \\
                \hline \hline
                     & bits   & az & az & as & az & as & az & az & - \\ \cline{2-10}
                     & uint   & az & az & as & cz & cs & az & - & az \\ \cline{2-10}
                     & int    & as & -  & as & -  & cs & -  & - &  - \\ \cline{2-10}
                from & ufixed & az & bz & bs & dz & dz & -  & - &  - \\ \cline{2-10}
                     & fixed  & as & -  & bs & -  & ds & -  & - &  - \\ \cline{2-10}
                     & log    & az & az & -  & -  & -  & a  & - &  - \\ \cline{2-10}
                     & float  & az & -  & -  & -  & -  & -  & e &  - \\ \cline{2-10}
                     & enum   & -  & az & -  & -  & -  & -  & - &  - \\
                \hline
                \end{tabular}

                \blist
                \item[{\bf a}] copy right justified
                \item[{\bf b}] copy from left of binary point, ignoring fractional part
                \item[{\bf c}] copy to left of binary point and zero fractional part
                \item[{\bf d}] copy aligned with binary point truncating or zero padding least significant
                  bits if size of fraction differs
                \item[{\bf e}] copy with adjustments to biased exponent and mantissa as required
                \blend

                \blist
                \item[{\bf z}] zero pad or truncate most significant end if width differs
                \item[{\bf s}] sign extend or truncate most significant end if width differs
                \blend

                If casting to or from compound types any types are allowed
                (including primitive types) as long as the numbers of bits are
                identical. Conversions into or out of compound types is carried
                out by mapping the bits in order without regard to component
                boundaries.

                When casting between int or uint and fixed or ufixed the binary
                point is recognised and where appropriate the sign is
                propagated, i.e. the cast retains numeric sense.

                When casting between bits and fixed or ufixed the binary point
                is not recognised and bits are simply right justified.

                When casting between bits and int or fixed and the cast type is
                wider, sign extension is performed instead of zero padding.

                Type float can only be cast to another float, adjusting the
                biased exponent and mantissa as required.

                Type bits or uint can be cast to type log which will extract the
                rightmost bit as false (0) or true (1). Casting type log to bits
                or uint results in a single bit, zero padded where required.

                Casting an enum to a uint provides the ordinal value (there is
                also an ord() built in function). Casting a uint to an enum
                selects the enum value given the ordinal.

                The function call is
                \begin{lstlisting}[gobble=16, language=threepl]
                   cast(expression, type)
                \end{lstlisting}
                where {\em type} is a string type specification or a type.


        \subsection{pointers}\label{sec:pointers}

            \index{type!pointer}
            \index{pointer}
            Pointers are immediate variables which can point to immediate or target
            variables, or components of compound variables. They can also point to
            classes, modules, procedures and functions. Target variables cannot
            themselves be pointers. Pointers can only be used during immediate
            execution, not target execution, i.e. all pointers are resolved before
            target code is generated. A pointer type is indicated by a {\em \verb'->'}
            in the type string. The type may be an array of pointers, e.g.
            \index{pointer!creation}
            \begin{lstlisting}[gobble=12, language=threepl]
               var(ptr1, "->");
               var(ptr2, "[10]->");
               var(ptr3, immediate, "[10]->");
            \end{lstlisting}
            The first example is a single pointer, the second an array of
            pointers. The 3rd example shows the redundant use of the {\em
            immediate} keyword - unnecessary since omission of the variable
            mode argument means that the variable is immediate.

            An unassigned pointer has a null value, which can be tested.

            Pointers may point to compound types (arrays, structs, subarrays,
            struct fields etc.).

            struct fields may be pointer types.

            Note that the type of the variable pointed to is not specified. A
            pointer may point to any type of variable and type checking is
            done when the pointer is resolved.

            \index{operator!\&}
            \index{type!pointer!assignment}
            \index{pointer!assignment}
            A pointer variable is assigned using the immediate assignment
            statement syntax
            \begin{lstlisting}[gobble=12, language=threepl]
               ptr1 = &v1;
               ptr2[i] = &v2;
               ptr3 = &println;
            \end{lstlisting}
            The {\em \verb'&'} operator gets a pointer to a variable, class, module,
            procedure or function. A pointer to a variable can be to the whole
            variable, a sub array, an array member or a struct field. A pointer can
            also point to a list or a map. However, a pointer cannot point to a list
            entry or a map entry or anything inside such an entry.
            
            A pointer reference from the {\em \verb'&'} operator can be -
            \blist
            \item   assigned to a pointer variable, array member or struct
                    field
            \item   passed as an argument to a pointer parameter
            \item   used as an initialiser to a created pointer variable
            \item   used as an initialiser within a compound value
            \item   appear within a compound value on the RHS of an
                    assignment statement
            \item   appear within a compound value which is an
                    argument in a module, procedure or function call
            \blend
            
            For modules, procedures and functions these can be either inbuilt or
            user-defined. Since variables and inbuilt and user-defined modules,
            procedures and functions have separate identifier tables, and across these
            tables identifiers may not be exclusive, the address operator {\bf
            \verb'&'} searches in the following order -
            \blist
            \item   variables
            \item   user-defined classes (there are no inbuilt classes)
            \item   inbuilt modules
            \item   user-defined modules
            \item   inbuilt procedures
            \item   user-defined procedures
            \item   inbuilt functions
            \item   user-defined functions
            \blend
            and returns the first that it encounters.

            \index{type!pointer!null assignment}
            \index{pointer!null assignment}
            A pointer may be explicitly assigned null by -
            \begin{lstlisting}[gobble=12, language=threepl]
               ptr = null;
            \end{lstlisting}

            \index{type!pointer!testing}
            \index{pointer!testing}
            Pointers may be tested for equality using the {\em ==} and {\em
            !=} operators, e.g. -
            \begin{lstlisting}[gobble=12, language=threepl]
               ptr1 == ptr2
               ptr3 != null
               ptr4 == &var[1].f
            \end{lstlisting}
            the other relational operators and arithmetic operators cannot be
            used on pointers.

            \index{type!pointer!resolving}
            \index{pointer!resolving}
            A pointer is resolved by using the operators {\em \verb'->'}. If
            the pointer points to a compound type then subscripts or field
            identifiers follow the {\em \verb'->'}. Subscripts and fields
            associated with the pointer variable are placed following the
            variable identifier and prior to the {\em \verb'->'}.
            
            A second version of the {\em \verb'->'} pointer operator can be used
            to dereference pointers in arguments to modules, procedures or
            functions. This operator is {\em \verb'->>'}. An attempt to
            dereference a pointer which is null using {\em \verb'->'} results
            in a fatal error. A null pointer dereferenced using {\em \verb'->>'}
            results in a null value. As an argument this is equivalent to
            omitting the argument. Thus if a module, procedure or function is
            called using arguments which are pointers dereferenced by using
            {\em \verb'->>'} then any pointer which is null will simply
            omit the argument. Use of {\em \verb'->'} in this case would
            give a fatal error. The two forms of pointer dereferencing are
            otherwise identical but the resulting null from the second form
            is not acceptable in other contexts so it would confer no advantage.

            The following example shows a single pointer to a single variable
            \begin{lstlisting}[gobble=12, language=threepl]
               var(ptr, "->");
               static(v1, "uint:8");
               static(v2, "uint:8");

               ptr = &v1;

               ptr-> <- 3;     // equivalent to v1 <- 3
               v2 <- ptr->;    // equivalent to v2 <- v1
            \end{lstlisting}
            Note that a pointer operator can be used on the left hand side of
            an assignment statement.

            For a pointer to an array
            \begin{lstlisting}[gobble=12, language=threepl]
               var(ptr, "->");
               static(v1, "[4]uint:8");
               static(v2, "uint:8");
               int(i, 1);

               ptr = &v1[0];

               ptr->[1] <- 3;     // equivalent to v1[1] <- 3
               v2 <- ptr->[i];    // equivalent to v2 <- v1[i]
            \end{lstlisting}

            For a pointer to a struct
            \begin{lstlisting}[gobble=12, language=threepl]
               struct(s,
                       f1,   "log",
                       f2,   "uint:8"
               );
               static(v1, s);
               static(v2, "log");
               var(ptr, "->");

               ptr = &v1;

               ptr->.f2 <- 3;     // equivalent to v1.f2 <- 3
               v2 <- ptr->.f2;    // equivalent to v2 <- v1.f2
            \end{lstlisting}

            For an array of pointers
            \begin{lstlisting}[gobble=12, language=threepl]
               struct(s,
                       f1,   "log",
                       f2,   "uint:8"
               );
               var(ptr, "[4]->");
               static(v1, "[4]s");
               static(v2, "uint:8");
               int(i, 3);

               ptr[0] = &v1[0].f1;
               ptr[1] = &v1[1];
               ptr[2] = &v1;

               ptr[0]-> <- false;
               v2 <- ptr[1]->.f2;
               when (ptr[2]->[i].f1)
                   ........;
            \end{lstlisting}

            For pointer to pointers to pointers etc. the pointer
            dereferencing operator is simply repeated, e.g.
            \begin{lstlisting}[gobble=12, language=threepl]
               p1[1]->.f1->[i].f2->
               p2->->->
            \end{lstlisting}

            Note that in the above examples where target variables are
            referenced using an immediate variable subscript, the value of
            the subscript at that time is substituted and thus the subscript
            in the target code is a constant. Target subscripts containing
            static and queue mode variables can be used for target variables
            and this is described later.
            
            \index{pointer!module}
            \index{pointer!procedure}
            \index{indirect!module call}
            \index{indirect!procedure call}
            Pointers to modules and procedures are dereferenced again using
            the {\bf \verb'->'} operator -
            \begin{lstlisting}[gobble=12, language=threepl]
               var(p, "->");
               .
               .
               p = &buffer;
               .
               .
               p->(p1 <- p2);
            \end{lstlisting}
            A pointer variable {\bf p} has been assigned a pointer to the
            inbuilt module {\bf buffer()} and then the module has been
            called via the pointer. Note that we could also have initialised
            the pointer variable -
            \index{type!pointer!initialisation}
            \index{pointer!initialisation}
            \begin{lstlisting}[gobble=12, language=threepl]
               var(p, "->", &buffer);
               .
               .
               p->(p1 <- p2);
            \end{lstlisting}
            Procedures are handled using the same syntax.
            
            \index{indirect!function call}
            \index{pointer!function}
            Pointers to functions are handled in a similar manner -
            \begin{lstlisting}[gobble=12, language=threepl]
               var(p, "->", &exists);
               .
               .
               if (p->(v))
                   println("variable '" + v + "' exists");
            \end{lstlisting}
            
            Where such a pointer is within an array or a struct subscripts
            and fields can be interposed -
            \begin{lstlisting}[gobble=12, language=threepl]
               var(p, "[3]->", {, &exists});
               .
               .
               if (p[1]->(v))
                   println("variable '" + v + "' exists");
            \end{lstlisting}
            
            Further indirection could be used -
            \begin{lstlisting}[gobble=12, language=threepl]
               var(p, "[3]->", {, &exists});
               var(q, "[4]str", {,, "p"});
               .
               .
               if (q[2]->p[1]->(v))
                   println("variable '" + v + "' exists");
            \end{lstlisting}
            
            Where a function returns a pointer the pointer can be
            dereferenced by applying the {\bf \verb'->'} operator to the
            function call, e.g.
            \begin{lstlisting}[gobble=12, language=threepl]
               c = f(a, b)->;
            \end{lstlisting}
            Subscripts and fields can be appear both before and after the
            \verb'->' where appropriate.
            




    \section{Dimensionality of variable references}

        \index{dimensionality}
        When assigning variables or passing variables as arguments to
        modules, procedures or functions compound types are allowed. Thus an
        array, sub-array or struct, or any legitimate combination of these,
        can be assigned or passed as an argument. Clearly when this is done
        the type structure must match across the assignment or
        argument/parameter match.

        A constant subscript reduces the dimensionality of an array. For
        example if we declare an array by -
        \begin{lstlisting}[gobble=8, language=threepl]
              var(a, "[4]int");
        \end{lstlisting}
        then {\bf a} is an array, {\bf a[1..2]} is an array, but {\bf a[1]} is not
        an array. Similarly for -
        \begin{lstlisting}[gobble=8, language=threepl]
           var(a, "[4][3][4]int");
        \end{lstlisting}
        {\bf a[][2][1..2]} is a 2-dimensional array since the constant
        subscript gets rid of one dimension. Trailing subscripts can be
        omitted without changing the result, so that {\bf a[1..2][2]} is a
        2-dimensional array as the omitted trailing subscript implies the
        full dimension.

        \index{assignment!arrays}
        \index{array!assignment}
        \index{assignment!structs}
        \index{struct!assignment}
        \subsection{compound assignment}

            Compound objects can be assigned as long as the dimensionality of the
            LHS and RHS match. For example -
            \begin{lstlisting}[gobble=12, language=threepl]
               var(a, "[4][3][4]int");
               var(b, "[6][4]int");
               a[1..2][1] = b[3..4];
            \end{lstlisting}
            is valid as each side has dimensions {\bf [2][4]}.

            Structs can be nested within types but the same dimensionality
            consideration applies -
            \begin{lstlisting}[gobble=12, language=threepl]
               struct(s,
                   a,  "[2][8]int",
                   b,  log
               );
               var(a, "[4][3][4]s");
               var(b, "[6][4]int");
               a[1..2][1][2].a[1][2..5] = b[3..4];
            \end{lstlisting}
            here each side again has dimensions {\bf [2][4]}.
            
            If the RHS is a compound value with missing elements, LHS
            elements corresponding to those missing on the right will
            not be changed.

        \subsection{compound arguments}\label{sec:compoundarg}

            \index{argument!compound}
            \index{compound!argument}
            When a compound argument is passed to a module, procedure or
            function parameter which has no type, the parameter takes
            on a type which matches the dimensionality of the argument.
            Any subscripts
            and field identifiers attached to the parameter apply to the
            compound object passed to the parameter. Thus in the
            following example -
            \begin{lstlisting}[gobble=12, language=threepl]
               var(a, "[4][3][4]int");

               .
               proc(a[1..2][1]);
               .

               procedure
               proc(var(p)) {
                   ...p...             // 2-dimensional array
                   ...p[0]...          // 1-dimensional array
                   ...p[0][1..2]...    // 1-dimensional array
                   ...p[][0]...        // 1-dimensional array
                   ...p[0][1]...       // single int
               }
            \end{lstlisting}
            the procedure parameter is passed a 2-dimensional array. The
            matching parameter is now handled in the procedure as though
            declared as a 2-dimensional array with dimensions appropriate to
            the array passed. As with assignments the parameter can be
            referenced as a 2-dimensional array, a 1-dimensional array (one
            subscript given) or a single element (two subscripts given).

            Again struct field identifiers may appear among the subscripts if
            structs are involved.
            
            If a module, procedure or function parameter is created as
            a compound type then the type of the argument must match
            that of the parameter.
            
            Between these extremes if a parameter contains undimensioned arrays
            then these will take on the dimensionality of the argument, but
            otherwise the types must match. If a parameter contains target types
            whose width specifiers are missing these will be supplied from the
            argument, but again the types must otherwise match. An undimensioned
            array parameter can derive missing array dimensions from the argument
            supplied or else from an initialiser if present. The argument and the
            initialiser may be variables or compound values. If a parameter
            containing undimensioned arrays is not initialised and is not passed
            an argument then the undimensioned arrays will be given dimensions
            0.
            
            \index{array!undimensioned}
            \index{dimensions!missing}
            If a parameter creates an undimensioned array because it is
            not passed an argument and has no initialiser then an execution
            error occurs. If a parameter has an undimensioned array type
            and the argument is optional then some initialisation must be
            supplied, e.g.
            \begin{lstlisting}[gobble=12, language=threepl]
               procedure
               p1 (var(p, "[]int", {})) {
                   ...
                   ...
               }
            \end{lstlisting}
            which will create an integer array of dimension 1. Since the
            initialisation entry is null no assignment is made and the
            single array member will have the default value, which in the case
            of an integer is 0.
            
            Note that if a list of identifiers is supplied as the first argument to
            {\bf var()}, rather than a single identifier, then missing dimensions in
            an array type cannot be inferred from arguments since all variables are
            constructed to be the same single type. However, missing dimensions may
            still be derived from an initialiser.
            
            A value mode parameter may be passed an immediate agument. If the
            argument is a compound type then the parameter must have a target type
            equivalent to the immediate type type of the argument. Arithmetic
            values will be checked to ensure that they fit within the size of
            the matching target sub-types. For example -
            
            \begin{lstlisting}[gobble=12, language=threepl]
               int(i, 4);  
               struct (st,
                   ls, "[2]log",
                   fi, "int:8"
               );

               p({{false, true}, i+30});

               procedure p (value(v, st))) {
                   .
                   .
               }
            \end{lstlisting}


    \section{Arrays of queues or memories}

        \index{array!queues}
        \index{array!memories}
        Arrays of queues and arrays of memories are not supported in 3PL. However,
        when these are required there is a way to effectively create and access
        these. Queues and memories can be created within a function or procedure
        as local variables and pointers to them returned in a pointer array. For
        example to create a one-dimensional array of queues the following code
        would suffice (this is a simplified version as the library procedure of
        the same name in {\em stdlib.3pl} will handle any number of dimensions) -
        \begin{lstlisting}[gobble=8, language=threepl]
           procedure
           queuearray (var(ptr) <- str(type)) {
               for (int(i) ; i<dimension(ptr) ; i++) {
                   queue(p, type);
                   ptr[i] = &p;
               }
           }
        \end{lstlisting}
        The output parameter is a 1-dimensional array of pointers. The
        procedure creates as many queues as there are pointers in that array.
        The queues are created inside the {\em for} loop which has an
        immediate statement block. Each queue has the same identifier but they
        are unique as each iteration of the loop has a new scope. A pointer
        to each queue is assigned to a member of the pointer array after it is
        created (variables survive the termination of their scope but cannot
        be accessed directly via an identifier, only through a pointer).

        Note also that the {\em for} loop has a statement block in brackets,
        which has its own scope. A single unbracketed variable creation
        procedure call in a {\em for} loop would create the variable in the
        scope containing the loop, so the following code - 
        \begin{lstlisting}[gobble=8, language=threepl]
           for (int(i) ; i<dimension(ptr) ; i++)
               queue(p, type);
        \end{lstlisting}
        would produce a multiply-defined variable (the code would be useless
        anyway since we have not assigned pointers to the queues).

        The queues in the array can now be used by the following constructions -
        \begin{lstlisting}[gobble=8, language=threepl]
           .
           .
           p[1]-> <<- ....;
           .
           .... <- p[i]->;
           .
           .
        \end{lstlisting}
        If the queues are compound (arrays or structs) then the subscripts and
        field identifiers follow the right arrow.




    \section{Directives}\label{sec:directives}
    
        \index{directive}
        Directives are intended as a means by which immediate program execution
        can control the compiler. Thus far 40 directives are recognised by
        the compiler interpreter (see below) but more may be implemented in the
        future. Since directives can be created and evaluated by the application
        code as well as by the compiler interpreter they can be used in a
        limited way as an extra variable mechanism.

        A directive is an entry in a global map. Each directive has a key and a
        value. The directive is stored in the map using the key to access the
        value. The key is a string. The value may be one of the 3PL immediate
        primitive types {\bf int}, {\bf log}, {\bf float} or {\bf str} (but not
        an enumerated type).

        Directives are entered by using the inbuilt procedure {\bf
        directive()} (\autorefph{sec:ipdirective}) and are extracted by
        the inbuilt function {\bf directive()} (\autorefph{sec:ifdirective}).
        The inbuilt function {\bf direxists()} (\autorefph{sec:ifdirexists}) can
        be used to check if a particular directive is present. The directive
        map can be obtained using the inbuilt function {\bf directives()}
        (\autorefph{sec:ifdirectives}).

        Directives may also be entered by using the -D command line option when
        calling 3PL (\autorefph{options}).

        A directive can also be entered during parsing (prior to immediate
        execution) by the inbuilt procedure {\bf presetdirective()}
        (\autorefph{sec:ippresetdirective}).

        Directives are global in scope.
        
        The following directives are recognised by the compiler or set by the
        compiler itself -
        
        \blist
        \item   "{\bf languageVersion}"
	\item   "{\bf compileOnly}"
        \item   "{\bf compilerMakeDate}"
        \item   "{\bf author}"
        \item   "{\bf ALU}"
        \item   "{\bf continuous}"
        \item   "{\bf currentDirectory}"
        \item   "{\bf date}"
        \item   "{\bf designName}"
        \item   "{\bf family}"
        \item   "{\bf FIFO}"
        \item   "{\bf filePath}"
	\item   "{\bf forceExec}"
        \item   "{\bf gatesNotLUTs}"
	\item   "{\bf intermediateOnly}"
        \item   "{\bf intTruncWarning}"
	\item   "{\bf listNets}"
	\item   "{\bf listTDEs}"
	\item   "{\bf listTDEsSel}"
	\item   "{\bf locations}"
        \item   "{\bf netlistcomment}"
        \item   "{\bf netlistDisplay}"
        \item   "{\bf netlistFile}"
        \item   "{\bf optimiseConnect}"
        \item   "{\bf outputDirectory}"
        \item   "{\bf OS}"
        \item   "{\bf parentDirectory}"
        \item   "{\bf part}"
        \item   "{\bf postProcessTool}"
        \item   "{\bf QUEUEREG}"
        \item   "{\bf reportFile}"
        \item   "{\bf rptToFile}"
        \item   "{\bf shell}"
        \item   "{\bf sigList}"
        \item   "{\bf skipPostProc}"
        \item   "{\bf sourceFile}"
        \item   "{\bf sourceVersion}"
        \item   "{\bf srlAddrWidth}"
        \item   "{\bf unass\_out}"
        \item   "{\bf verbose}"
        \blend

        The compiler itself creates the following directives -
        \index{directive!arch}
        \index{directive!OS}
        \index{directive!cellName}
        \index{directive!compilerMakeDate}
        \index{directive!currentDirectory}
        \index{directive!date}
        \index{directive!designName}
        \index{directive!family}
        \index{directive!languageVersion}
        \index{directive!netlistFile}
        \index{directive!outputDirectory}
        \index{directive!parentDirectory}
        \index{directive!sourceFile}
        \index{directive!sourceVersion}
        \index{directive!srlAddrWidth}

        \begin{tabular}{| l | l | p{10.0cm} |}
        \hline
        "languageVersion"  & str & the compiler language version
        \\ \hline
        "sourceVersion"    & str & the current compiler source code version
                                   number
        \\ \hline
        "compilerMakeDate" & str & the date the compiler was made. Note that the
                                   date and time that the 3PL compiler is
                                   executed can be obtained using inbuilt
                                   functions \autorefph{sec:iftime} {\bf
                                   time()} for date and time elements and
                                   \autorefph{sec:iftimestring} {\bf
                                   timestring()} for date and time as a
                                   string.
        \\ \hline
        "OS"               & str & the host computer operating system
        \\ \hline
        "arch"             & str & the host computer architecture
        \\ \hline
        "srlAddrWidth"     & int & the address width of CLB shift registers
        \\ \hline
        "date"             & str & the current date and time
        \\ \hline
        "sourceFile"       & str & the source file name specified on the
                                   command line (with ".3pl" appended by 3PL if
                                   absent)
        \\ \hline
        "parentDirectory"  & str & the full path name of the directory containing the
                                   main source file
        \\ \hline
        "currentDirectory" & str & the full path name of the current directory
        \\ \hline
        "outputDirectory"  & str & the full path name of the output directory specified by the -o
                                   command line option, or the parent
                                   directory by default
        \\ \hline
        "designName"       & str & the design name derived from the source
                                   file name or specified by the -N command
                                   line option
        \\ \hline
        "cellName"         & str & the cell name to use if the netlist is a logic
                                   core produced by calling inbuilt procedure
                                   {\bf export()} \autorefph{sec:ipexport}.
        \\ \hline
        "family"           & str & a string designating the FPGA family, e.g. XC5V
%        \\ \hline
%        "simulation"             & log & true if the -s command line option has
%                                         been used - the simulator will be
%                                         executed following immediate execution
        \\ \hline
        \end{tabular}

        The compiler creates the following directives from command
        line options. Most of these directives can also be set by application code
        which will override any value derived from command line options.
        \index{directive!compileOnly}
        \index{directive!forceExec}
        \index{directive!gatesNotLUTs}
        \index{directive!intTruncWarning}
        \index{directive!intermediateOnly}      
        \index{directive!listNets}
        \index{directive!listTDEsSel}
        \index{directive!listTDEs}
        \index{directive!locations}
        \index{directive!netlistDisplay}
        \index{directive!optimiseConnect}
        \index{directive!rptToFile}
        \index{directive!sigList}
        \index{directive!skipPostProc}
        \index{directive!verbose}

        \begin{tabular}{| l | l | p{10.0cm} |}
        \hline
        "compileOnly"       & log &  true if the -c command line option has
                                     been used. 3PL will exit after compiling
                                     the source code and before execution.
                                     Cannot be set during execution.
        \\ \hline
        "forceExec"         & log &  true if the -f command line option has
                                     been used. If true immediate execution
                                     and netlist code generation will be
                                     carried out regardless of the time of
                                     last modification of associated files.
                                     Cannot be set during execution.
        \\ \hline
        "gatesNotLUTs"      & log &  true if the -G command line option has
                                     been used. If true this will generate
                                     netlist logic encoded as gates (AND, OR
                                     etc.) rather than as LUTs.
        \\ \hline
        "intermediateOnly"  & log &  true if the -i command line option has
                                     been used. 3PL will exit after compiling
                                     the source code and completing execution
                                     but prior to netlist generation. It can be
                                     invoked if the aim is just to examine the
                                     TDE list rather than to produce a netlist.
        \\ \hline
        "intTruncWarning"   & log &  true if the -w command line option has
                                     been used. 3PL will emit a warning to the
                                     report file for any target "int" or "uint"
                                     that has been assigned to a narrower
                                     destination. "fixed" and "ufixed" warnings
                                     are always given. If set during execution
                                     it wil take effect from then on. Similarly
                                     can be cleared during execution.
        \\ \hline
        "listNets"          & log &  true if the -n command line option has
                                     been used. This will produce a list
                                     of nets used in the netlist file with
                                     all aliases. The report is written to
                                     a .net file.
        \\ \hline
        "listTDEs"          & log &  true if the -t or -T command line options
                                     have been used. Either of these options
                                     may optionally be followed by the digit
                                     0, 1, 2 or 3 which will print the TDE list -\\
                                     & & 0 - after all optimisation \\
                                     & & 1 - before any optimisation \\
                                     & & 2 - after the 1st optimisation \\
                                     & & 3 - after family-specific optimisation. \\
                                     & & If no digit is encountered then the TDE list
                                     after all optimisation is printed.
        \\ \hline
        "listTDEsSel"       & int &  this directive contains an integer value
                                     derived from the optional digit described
                                     in "listTDEs" above. It is 0 by default.
        \\ \hline
        "locations"         & log &  true if the -l command line option has
                                     been used. If true listings of the TDE
                                     list will contain source code locations
        \\ \hline
        "netlistDisplay"    & log &  true if the -d command line option has
                                     been used. If true the interactive
                                     netlist display window will pop up
                                     following netlist generation.
        \\ \hline
        "optimiseConnect"   & log &  false if the -O command line option has
                                     been used; true by default. This disables
                                     code which attempts to eliminate CONNECT
                                     TDEs and connect them through the catalog
                                     instead. It can be used to diagnose some
                                     logic inconsistencies should they arise. It
                                     is only used during development or
                                     maintenance.
        \\ \hline
        "rptToFile"         & log &  true if the -r command line option has
                                     been used - verbose mode is enabled and
                                     the report output is written to file
                                     {\bf {\em design\_name}.rpt}
        \\ \hline
        "sigList"           & log &  true if the -T command line option has
                                     been used. It causes the TDE signal definition
                                     section of the TDE list to be printed along
                                     with the list itself. This is usually not
                                     usefull.
                                     {\bf {\em design\_name}.rpt}
        \\ \hline
        "skipPostProc"      & log &  true if the -s command line option has
                                     been used. It causes any post-processing
                                     3PL code to be skipped.
        \\ \hline
        "verbose"           & log &  true if the -v, -r, -t, -T or -P command
                                     line options have been used - a detailed
                                     report will be written to standard output by
                                     default or to file {\bf {\em
                                     design\_name}.rpt} if the -r option has been
                                     given. If set during execution some initial
                                     report output will be missing.
        \\ \hline
        \end{tabular}


        The remaining directives used by 3PL are set by application code -

        \index{directive!author}
        String directive "{\bf author}" is optional. If given a value
        (string) this will be inserted as an author field in the EDIF file
        (Xilinx claim this field should contain the company name, though there
        seems no reason why it matters what it contains). The compiler
        interpreter uses the value of the directive (if any) at the completion
        of immediate code execution.

        \index{directive!ALU}
        String directive "{\bf ALU}" specifies that any static or
        selectvalue variables which are directly assigned the result of an
        arithmetic operation\footnote{At the moment only + and -.} in two or
        more statements will have those operations combined into a single
        Arithmetic Logic Unit (ALU). Where a value to be added, subtracted or
        assigned is a constant, such constants are optimised to only appear
        once. This optimisation only applies to an operator whose result is
        directly assigned to the variable, not to operators in sub-expressions
        since these are evaluated in parallel and must have separate operator
        logic.
        
        Where a number of increments, decrements, additions or subtractions
        are made to a static variable the use of an ALU avoids duplication of
        expressions for each instance using instead a single combined
        adder/subtracter. In simple cases this substantially reduces
        the use of logic gates and does not increase path delays. However
        if any values to be assigned, added to or subtracted from the variable
        require combinatorial logic of significant depth, i.e. they involve
        further arithmetic operators, then the use of a shared ALU may
        result in greater path delays which may exceed the clock domain
        timing constraint.
        
        The effect of this directive can be changed for individual
        static variables by setting the {\bf ALU} attribute. If the directive
        is not defined or is not true then the ALU attribute can be set
        true on individual static variables to implement an ALU if
        appropriate (as per the mixed assignment types described above).
        If the directive is defined and is true setting attribute ALU
        false will disable the effect on a particular static variable.

        \index{directive!continuous}
        String directive "{\bf continuous}" is defined and true (log)
        the following logic simplifications will be applied -
        \blist
        \item   conditional static and queue assignments (within a {\bf when}
                statement) will be enabled directly from the {\bf when} test
                expression rather than indirectly via an execution signal
                through a {\bf when} TDE. This optimisation occurs if the
                directive exists and is true on encountering a when during
                immediate execution.
        \item   during optimisation following immediate execution all
                unconditional static and queue assignments will be removed from
                execution threads and enabled continuously. This achieves the
                same effect as setting the {\bf continuous} attribute true on
                each static variable. This optimisation occurs if the
                directive exists and is true at completion of immediate execution.
        \blend
        To enable both optimisations universally "{\bf continuous}"
        should be created and set true. This should be done using
        presetdirective() rather than directive() since it should be set
        prior to immediate execution - if set during immediate execution prior
        code may be generated without the directive which might introduce some
        subtle problems. The use of directive() for this purpose is not
        prohibited since it may be required to be set conditionally, in which
        case the conditional code should precede any  target code. Use of this
        directive can significantly reduce logic path delays and to a small
        extent the logic resources used.

        Directive "{\bf continuous}" can also be set by supplying the 3PL
        command line argument -D continuous.
        
        To restrict this optimisation to individual variables the "continuous"
        attribute can be set true on a static or queue mode variable instead
        of using the global directive. Alternatively the global effect can
        be removed for individual static or queue mode variables by setting
        the "continuous" attribute to false. The directive prevails if the
        attribute has not been set, otherwise the attribute is used.

        \index{directive!FIFO}
        Directive "{\bf FIFO}" when set true prior to completion
        of immediate execution specifies that queue mode variables whose depth is
        such that block RAMs will be used (generally depth \verb'>' 16) will be
        implemented using FIFO hardware associated with the block RAMs if the
        FPGA includes such hardware. This feature can be controlled on
        individual queue mode variables by attribute {\bf FIFO} which if
        specified overrides the directive. The default behaviour if neither
        directive nor attribute are used is to NOT use the dedicated FIFO
        hardware.
        
        Use of the FIFO hardware is a compromise between speed and logic
        usage. The FIFO hardware saves CLB gates and flip-flops, which would
        otherwise be used, but has greater path delays which may cause
        failure to meet timing constraints.
        
        \index{directive!filePath}
        Directive "{\bf filePath}" when true causes error messages to
        give file paths rather than file names. This may occasionally
        be useful to disambiguate variables of the same name in source
        files in different directories. This can also be set by the command
        line option -F.

        \index{directive!gatesNotLUTs}
        String directive "{\bf gatesNotLUTs}" when set true prior to
        completion of immediate execution causes many LUTs normally produced by
        the code generated to be replaced with nested gates. The logic remains
        the same but the use of nested gates instead of explicit LUTs may change
        the way the vendor place-and-route software optimises and packs the
        logic which in turn may change the ease, or lack thereof, with which the
        timing contraints can be met. Thus when timing constraints fail to be
        met, trying the -G option may improve the situation. The directive can
        also be set using the -G option when calling 3PL

        \index{directive!netlistcomment}
        String directive "{\bf netlistcomment}", if defined, is copied as
        a comment to the EDIF netlist file and the NCF contraints file or the
        XDC constraints file. In the EDIF netlist file it appears within double
        quotes in a comment field. In the constraint files it appears following
        a "\#". If the directive contains more than one line these will appear
        as separate comment lines in the constraint files, but in the netlist
        file the whole string will appear within a single pair of quotes,
        including the newlines.

        \index{directive!part}
        String directive "{\bf part}"must contain the full FPGA part
        specification (string). It is set within inbuilt procedure
        {\bf part()}(\autorefph{sec:ippart}). The part specification is not only inserted in
        the EDIF netlist file, it is also used to guide code generation to suit
        the FPGA type. The directive must be assigned a value prior to target
        code generation and its value cannot be changed during  the remainder of
        immediate execution. An attempt to change the directive will result in a
        warning message. If this directive is not recognised a fatal error will
        be flagged and compilation will terminate. If this directive is not
        defined, at the moment (with single vendor implementation) the Xilinx
        UNISYMS generic part will be used.
        
        Allowed values for directive "part" are -                       
        \blist
        \item   "XC95..." XC9500 CPLD, e.g. "XC9536PC44-5"
        \item   "XC2C..." Coolrunner 2 CPLD, e.g. "XC2C32APC44-6"
        \item   "XC2S..." Spartan 2 FPGA
        \item   "XC3S..." Spartan 3 FPGA
        \item   "XC6S..." Spartan 6 FPGA
        \item   "XCV..." Virtex or Virtex E FPGA, e.g. "XCV300EBG432-7"
        \item   "XC2VP..." Virtex 2 Pro FPGA, e.g. "XC2VP40FF1152-6"
        \item   "XC2V..." Virtex 2 FPGA, e.g. "XC2V3000BF957-4"
        \item   "XC4V..." Virtex 4 FPGA
        \item   "XC5V..." Virtex 5 FPGA
        \item   "XC6V..." Virtex 6 FPGA
        \item   "XC7A..." Artix 7 FPGA
        \item   "XC7K..." Kintex 7 FPGA
        \item   "XC7V..." Virtex 7 FPGA
        \item   "XC7Z..." Zynq FPGA
        \blend

        \index{directive!QUEUEREG}
        String directive "{\bf QUEUEREG}" when set true prior to
        completion of immediate execution specifies that queue mode variables are
        to have a pipelined register added to their outputs. This reduces the
        clock-to-output time. The queue buffer logic incorporates the extra
        buffer with no change to functionality. The attribute {\bf QUEUEREG} may
        be set (true or false) on individual queue mode variables, overriding
        the global directive.

        \index{directive!reportFile}
        String directive "{\bf reportFile}" contains the absolute path name
        of the report file.

        \index{directive!shell}
        If the {\bf shell()} \autorefph{sec:ipshell} procedure is to be used and
        a shell other than the default shells {\bf /usr/bin/bash} or {\bf
        /bin/bash} is required then the string directive {\bf shell} should be
        defined with the value being a string giving the shell path.

        \index{directive!unass\_out}
        String directive "{\bf unass\_out}" controls the behaviour of
        output mode variables which have unassigned components. The directive
        type "log" -
        \blist
        \item   false - connect to ground
        \item   true - connect to VCC
        \blend
        For individual output mode variables the default can be controlled
        by the output mode variable attribute {\bf defaultout} (see
        \autorefph{sec:attributes}). The action of this attribute overrides
        directive "unass\_out". If neither specify a values an unassigned output
        mode variable component will result in a fatal error.

        \index{directive!postProcessTool}
        String directive "{\bf postProcessTool}" is mandatory. It should have
	the value "ise" (for the Xilinx ISE tools) or "vivado" (for Xilinx
	Vivado).  {\bf It must be defined
        during parsing by presetdirective() or fatal termination of 3PL
        immediate execution will occur.} This should be done in the header file
        for each FPGA family member. If the directive is changed during
        immediate execution it must be remembered that post-processing 3PL code
        will use this directive and the default value for the FPGA platform
        will prevail if post-processing is not preceded by main-code execution
        because the EDIF file is already up-to-date from a previous run.
        
        The choice of place-and-route tools changes 3PL in two ways -
        \begin{enumerate}
        \item   3PL will create either an NCF constraint file for ISE or an
                XDC constraint file for Vivado, assuming that constraints
                have been generated
        \item   the form of element pins in the EDIF netlist for arrays is
                different, ISE requiring pin arrays of the form A0, A1, A2,
                A3 and Vivado A[0:3]
        \end{enumerate}
        
        
        

        
        

    \section{Attributes}\label{sec:attributes}

        \index{attributes}
        Attributes are assigned to variables to 
        \blist
        \item   control some aspects of code generation
        \item   set associated parameters such as clock frequency or queue buffer depth
        \item   pass constraints to the EDIF netlist generator
        \blend
        
        Case is ignored in attribute name strings. They are converted to lower
        case when used within 3PL.
        
        Attributes are passed by means of a map type. This may be a variable
        of type "map" or it may be a one-dimensional compound value where every
        entry has the key=value format, the key being the attribute identifier
        and the value its associated value.
        
        Among the attributes described below the {\bf loc} attribute applies
        to several variable modes. This attribute is usually an array of
        strings ("[]str") but where a single location is involved it may also
        be a string ("str").

        Attributes are usually simply omitted if not required, but they
        may be given default or inactive values if desired.
        Where an attribute is boolean (type "log") it must be present and
        true to take effect.      
        
        Attributes are described within the variable creation procedures to
        which they apply.
        
        \subsection{applying attributes}
        
            Attributes are applied to a variable in a map. The map
            keys are the attribute identifiers and the associated values
            are the attribute values. The entry values are mixed types,
            the type depending on the individual attribute.
            
            The variable creation inbuilt procedures
            {\bf input()}(\autorefph{sec:ipinput}),
            {\bf output()}(\autorefph{sec:ipoutput}) and
            {\bf clock()}(\autorefph{sec:ipclock})
            have an optional final argument which is an attribute map.
            The attribute map argument can be a variable of type "map"
            or it can be a compound value. The following code
            segments all achieve the same thing -

            \begin{lstlisting}[gobble=12, language=threepl]
               var(cmap, "map");
               cmap["frequency"] = 100.0e6;
               cmap["dutycycle"] = 0.25;
               clock(c);
               attributes(c, cmap);
            \end{lstlisting}
            The {\bf attributes()}(\autorefph{sec:ipattributes}) inbuilt
            procedure applies an attributes map to a variable.
            This is very longwinded,
            \begin{lstlisting}[gobble=12, language=threepl]
               var(cmap, "map", {frequency=100.0e6, dutycycle=0.25});
               clock(c);
               attributes(c, cmap);
            \end{lstlisting}
            is still rather verbose,
            \begin{lstlisting}[gobble=12, language=threepl]
               var(cmap, "map", {frequency=100.0e6, dutycycle=0.25});
               clock(c, cmap);
            \end{lstlisting}
            is shorter.
            
            Attributes can instead be passed directly using key=value arguments,
            e.g.
            \begin{lstlisting}[gobble=12, language=threepl]
               clock(c, frequency=100.0e6, dutycycle=0.25);
            \end{lstlisting}
            is the most concise.
            
            The {\bf attributes()} inbuilt procedure will accept a bracketed
            list of variables as the first argument, applying the attributes
            to every variable in the list.
            
            The {\bf attributes()} procedure does not enter a new attributes map in
            the variable, it adds the entries in the argument map to that in the
            variable. Thus the attributes() procedure can be called multiple times
            to add one or more attributes to the variable. If the same attribute is
            entered again the new value replaces the old value. This can be used to
            enable and disable 3PL attributes that are used during immediate
            execution.
            
            Netlist attributes can be deleted by assigning the value {\bf null}.
            3PL attributes cannot be deleted; 3PL attributes either
            relate to parameters that are mandatory such as clock domains, or are
            boolean in which case they can be disabled rather than deleted, so
            there is no reason to delete them.

            Netlist attributes are only used at the termination of immediate
            execution so changing their values or deleting them during immediate
            execution is not an error.
            
            Attributes {\bf domain}, {\bf readdomain} and {\bf writedomain}
            can be assigned {\bf null} and this sets the associated
            domain to the current clock domain.
            
            Where a variable is passed as an argument to a module, procedure
            or function the associated parameter is implemented as an
            internal pointer to the argument variable (call by reference)
            and hence if attributes are changed in the parameter they are
            changed in the argument variable.
            
            

        \subsection{examining attributes}
            
            The attributes of a variable can be fetched by the
            inbuilt function {\bf attributes()}(\autorefph{sec:ifattributes})
            which returns a map. If the second optional argument to the
            function is true an extended set of attributes is returned. These
            are -
            \blist
            \item    {\bf identifier} - the abbreviated identifier
            \item    {\bf eidentifier} - the extended identifier
            \item    {\bf mode} - the mode
            \item    {\bf type} - the type
            \item    {\bf typestring} - the type as a string
            \item    {\bf decloc} - the source file location where created
            \item    {\bf inputparam} - true if this is an input parameter
            \item    {\bf outputparam} - true if this is an output parameter
            \item    {\bf matched} - true if this is a parameter matched by an argument
            \item    {\bf readonly} - true if read-only
            \item    {\bf assigned} - true if the variable has been assigned
            \item    {\bf used} - true if it is a target mode variable and its
                                value has been used
            \item    {\bf clk\_srcs} - for a clock variable the number of sources
            \item    {\bf clk\_sinks} - for a clock variable the number of sinks
            \item    {\bf clk\_CE} - for a clock variable the number of element
                                clock inputs driven by this clock. This may be
                                smaller than {\bf clk\_sinks} as that attribute
                                includes clock inputs to buffers or clock managers
                                where the clock is not triggering state changes.
            \blend
            The attributes in this list are read-only and cannot be assigned to
            the variable.
            
            A single attribute can be fetched by the inbuilt function
            {\bf attribute()}(\autorefph{sec:ifattribute}).
            
            Attributes can also be fetched for listing purposes as an array of
            structs by calling {\bf listattributes()} \autorefph{sec:iflistattrs}.
            The struct contains three string fields, key, the attribute key,
            mode, the mode of the value, and value which is a string version
            of the value or is the identifier of a clock domain. This array
            may be printed to standard output by calling library procedure
            {\bf printattributes()} \autorefph{sec:lpprintattrs}.
            
            Attribute {\bf domain} for a value mode variable has a
            null value when read back regardless of any assignments to
            the variable.

            For a selectvalue, static or priority mode variable attributes
            {\bf domain}, {\bf readdomain} and {\bf writedomain} are the same.
            These attributes are only added to the variable when it is
            read or written.
            
            For a queue mode variable, attributes {\bf readdomain}, {\bf
            writedomain} and {\bf domain} are the same if the queue is
            synchronous. If the queue is asynchronous attributes {\bf readdomain}
            and {\bf writedomain} differ and attribute {\bf domain} is not
            present in the attribute map. These attributes are only added to the
            variable when it is read ({\bf readdomain}), written ({\bf writedomain})
            or both ({\bf domain} if synchronous).
            
            Attribute {\bf identifier} is the variable identifier, i.e. the
            identifier used when the variable was created.

            Attribute {\bf eidentifier} is the expanded identifier. This is
            full unambiguous identifier used in the intermediate code. Both
            scope information and argbitrary appended integers are used to
            construct a unique identifier.
            
            Attribute {\bf mode} is the variable mode as a string.
            
            Attribute {\bf type} is the variable type returned as a "type" type.

            Attribute {\bf typestring} is the variable type returned as a string.

            Attribute {\bf decloc} is the source code location where the variable
            was created.
            
            Attribute {\bf inputparam} is a value of type "log" which is true
            if this variable is an input parameter to a module, procedure or
            function.
            
            Attribute {\bf outputparam} is a value of type "log" which is true
            if this variable is an output parameter to a module, procedure or
            function.
            
            Attribute {\bf matched} is a value of type "log" which is true
            if this variable is a parameter to a module, procedure or
            function and has been passed an argument.
            
            Attribute {\bf readonly} is a value of type "log" which is true
            if this variable has been made read-only.
            
            Attribute {\bf assigned} is a value of type "log" which is true if
            this variable has been assigned. The variable cannot be input mode,
            cmemory mode or rmemory mode.
            
            For an immediate and value mode variables this attribute is true if
            the variable has been explicitly assigned a value during immediate
            execution - this does not include any initial value given to the
            variable on creation.

            For target variables other than value mode this attribute is true
            if a target assignment statement has been processed during immediate
            execution where this variable appears on the left hand side - this
            excludes any initial values given to variables on creation.

            For clock mode variables this attribute is true if -
            \blist
            \item   the clock variable was created using inbuilt procedure {\bf
                    clock()} where external input has been specified by
                    supplying the location attribute
            \item   the clock variable is a negative clock that has been created
                    using inbuilt procedure negclock() and the positive clock from
                    which it is derived is later assigned or has previously been
                    assigned
            \item   the clock variable has been assigned from another clock variable
                    which is later assigned or has previously been assigned
            \item   the clock variable is an input parameter to inbuilt
                    procedure elementoutput() where it is connected to an
                    output port or a 3-state port
            \item   the clock variable is an output parameter to inbuilt
                    procedure clockfromexpr()
            \blend

            Attribute {\bf clk\_sinks} for a clock variable is the number of
            places the clock is used to clock logic

            Attribute {\bf clk\_srcs} for a clock variable is the number of
            places the clock is driven. If this is not 1 either the clock has no
            source or it has multiple sources; in either case an error.


           
       
        




    \section{Operators}

        \index{operator}
        The following operators have been implemented.

        \begin{tabular}{| l | l | l |}
        \hline
        \verb'~' & unary prefix & bitwise inversion \\
        \hline
        \verb'!' & unary prefix & logical negation \\
        \hline
        \verb'-' & unary prefix & arithmetic negation \\
        \hline
        \verb'<' & unary prefix & queue read availability \\
        \hline
        \verb'>' & unary prefix & queue write availability \\
        \hline
        \verb'&' & unary prefix & address \\
        \hline
        \verb'++' & unary prefix & increment \\
        \hline
        \verb'--' & unary prefix & decrement \\
        \hline
        \verb'++' & unary postfix & increment \\
        \hline
        \verb'--' & unary postfix & decrement \\
        \hline
        \verb'?' & unary prefix & queue examine \\
        \hline
        \verb'->' & unary postfix & pointer dereference \\
        \hline
        \verb'->>' & unary postfix & pointer dereference \\
        \hline
        \verb'*' & binary & multiplication \\
        \hline
        \verb'/' & binary & division \\
        \hline
        \verb'%' & binary & remainder \\
        \hline
        \verb'+' & binary & addition or string concatenation\\
        \hline
        \verb'-' & binary & subtraction \\
        \hline
        \verb'<<' & binary & left shift \\
        \hline
        \verb'>>' & binary & right shift \\
        \hline
        \verb'^<<' & binary & left rotate \\
        \hline
        \verb'>>^' & binary & right rotate \\
        \hline
        \verb'<' & binary & less than  \\
        \hline
        \verb'>' & binary & greater than \\
        \hline
        \verb'<=' & binary & less than or equal to \\
        \hline
        \verb'>=' & binary & greater than or equal to \\
        \hline
        \verb'==' & binary & equal to \\
        \hline
        \verb'!=' & binary & not equal to \\
        \hline
        \verb'&' & binary & bitwise AND \\
        \hline
        \verb'^' & binary & bitwise XOR \\
        \hline
        \verb'|' & binary & bitwise OR \\
        \hline
        \verb'&&' & binary & logical AND \\
        \hline
        \verb'||' & binary & logical OR \\
        \hline
        \verb'^^' & binary & logical XOR \\
        \hline
        \verb'+=' & binary & add or concatenate right to left \\
        \hline
        \verb'-=' & binary & subtract right from left \\
        \hline
        \verb'*=' & binary & multiply left by right \\
        \hline
        \verb'/=' & binary & divide left by right \\
        \hline
        \verb'%=' & binary & remainder of left divide by right \\
        \hline
        \verb'&=' & binary & bit AND right to right \\
        \hline
        \verb'|=' & binary & bit OR right to right \\
        \hline
        \verb'^=' & binary & bit XOR right to right \\
        \hline
        \verb'&&=' & binary & logical AND right to right \\
        \hline
        \verb'||=' & binary & logical OR right to right \\
        \hline
        \verb'^^=' & binary & logical XOR right to right \\
        \hline
        \verb'<<=' & binary & left shift left by right \\
        \hline
        \verb'>>=' & binary & right shift left by right \\
        \hline
        \verb'?:' & ternary & conditional \\
        \hline
        \end{tabular}
        
        Note that the C ',' operator (comma), as sometimes used in {\bf for} statements is
        not implemented in 3PL.




        \index{operator!precedence}
        The binary and unary operators are listed in order of decreasing
        precedence in the following table -

        \begin{tabular}{| l |}
        \hline
        \verb'*' \verb'/' \verb'%' \\ \hline
        \verb'+' \verb'-' \\ \hline
        \verb'==' \verb'!=' \verb'<' \verb'>' \verb'<=' \verb'>=' \\ \hline
        \verb'&' \\ \hline
        \verb'^' \\ \hline
        \verb'|' \\ \hline
        \verb'&&' \\ \hline
        \verb'^^' \\ \hline
        \verb'||' \\ \hline
        \verb'?:' \\ \hline
        \end{tabular}

        Note that the relational operators have higher precedence than the
        logical operators so that an expression such as \verb'a < 3 && b > 5'
        does not need parentheses.

        \index{operator!\verb'++' \verb'--'}
        The {\bf ++} and {\bf --} operators can be used on {\bf immediate}
        and {\bf static} variables, but not {\bf value}, {\bf selectvalue},
        {\bf queue} or {\bf priority} variables.
        
        For {\bf static} variables these operators can only be used in the form of
        increment and decrement statements and not in any other context, e.g. in an
        expression. Thus for example for static variable {\bf t} -
        \begin{lstlisting}[gobble=8, language=threepl]
           t++;
        \end{lstlisting}
        is allowed but the forms
        \begin{lstlisting}[gobble=8, language=threepl]
           a <- t++;       // ERROR!
           a[t++] <- b;    // ERROR!
           func(t++);      // ERROR!
        \end{lstlisting}
        are not and would result in a fatal immediate execution error.
        
        The pre-increment and post-increment forms behave identically since they
        form a statement and not part of an expression. The same applies to the
        decrements.

        \index{operator!unary \verb'?'}
        The {\bf \verb'?'} prefix unary operator allows evaluation of an
        expression containing queues without removing data from the queues.
        Unlike a queue read this operation proceeds regardless of queue
        availability. If an examined queue is not available the data is not
        valid. Use of this operator would normally be accompanied by an
        explicit queue availability test using the {\bf \verb'<'} prefix unary
        operator (next paragraph). The operand may be a single queue variable
        or any expression containing queues, but need not contain queues at all
        and may even be immediate. If the operand contains no queue reads then
        the value is simply returned.

        \index{operator!unary \verb'<'}
        \index{queue!availability}
        The {\bf \verb'<'} prefix unary operator examines its expression argument
        for queue read availability and produces a logical result. This operator does
        not execute a read, i.e. a datum will not be removed from any queues. The
        expression may be a single queue variable or may contain a number of queues,
        but the expression need not contain queues and may even be immediate.
        Where the operand contains no queues {\bf \verb'<'} will return the constant
        {\em true}. It would be usual where the argument is an expression rather
        than a single queue variable for the expression to be passed via a value
        variable since the expression is presumably used elsewhere for actual
        evaluation and duplicating the code and then ignoring the result by applying
        the operator would be silly. Queue references within the operand of the {\bf
        \verb'<'} operator are ignored by any enclosing {\em sync}. The read
        availability operator can only be used in a module within which one or more
        reads of the queue or queues are also executed (read availability of a queue in
        a module which does not read the queue is not meaningful). The resulting
        logical value has the output (read) clock domain of the queue variable. If
        the output (read) clock domain has not yet been established then use of the
        availability logical value in a clocked context will establish the output
        (read) clock domain of the queue variable\footnote{If the queue read
        availability operator is used in an expression assigned to a value mode
        variable the output clock domain of the associated queue variable will not be
        established by that assignment. If the value variable is subsequently used
        on the RHS of an assignment to a clocked variable, or in another context in
        which a clock domain is required, then that will establish the clock domain
        of the associated queue variable if it is not already established.}.

        \index{operator!unary \verb'>'}
        \index{queue!availability}
        The {\bf \verb'>'} prefix unary operator examines its variable argument for
        queue write availability and produces a logical result. This operator does
        not execute a write. Where the variable has one or more target subscripts
        the {\bf \verb'>'} operator requires them to exhibit read availability as
        well before it will return true. The resulting logical value has the input
        (write) clock domain of the queue variable. If
        the input (write) clock domain has not yet been established then use of the
        availability logical value in a clocked context will establish the input
        (write) clock domain of the queue variable.

        Note that the {\bf \verb'>'} operator requires the queue variable
        assignment clock domain. It must be called in the same clock domain as
        assignments to the queue and if its use precedes any assignments in the
        source code it will set the assignment clock domain.

        \index{operator!\verb'->'}
        \index{indirect}
        \index{pointer!dereference}
        The postfix {\bf \verb'->'} and {\bf \verb'->>'} unary operators
        dereference a pointer. The second form is for use in a module, procedure
        or function argument where a null pointer results in omission of the
        argument rather than a fatal error. Note that this differs from its use in
        C in that if a struct is pointed to a field reference still requires a "."
        before the field identifier, e.g.
        \begin{lstlisting}[gobble=8, language=threepl]
           p->.field
        \end{lstlisting}
        \index{indirect from string}
        \index{pointer!dereference from string}

        The {\bf \verb'->'} unary operator will also dereference a string variable
        where the string is a variable identifier. For example
        \begin{lstlisting}[gobble=8, language=threepl]
           str(s, "var1");
           println(s->);
        \end{lstlisting}
        will print the value of variable {\bf var1}. The identifier string
        may include scope prefices such as "global.".
        
        The {\bf \verb'->'} operator may be preceded by fields or subscripts
        to dereference a pointer or string contained in a compound type.

        The {\bf \verb'->'} operator may be followed by fields or
        subscripts when the type pointed to is compound.

        The {\bf \verb'->'} unary operator will also dereference an expression
        in parentheses, e.g.
        \begin{lstlisting}[gobble=8, language=threepl]
           int(i, 5);
           println(("var" + i)->);
        \end{lstlisting}
        will print the value of variable {\bf var5}. The expression may include
        function calls and can be of type "str" or type \verb'"->"'. In the latter
        case the expression can only be a function returning a pointer since
        arithmetic cannot be performed on pointers. The expression could be
        the address of a variable using the \verb'&' operator but that would
        be singularly pointless.

        \index{operator!ternary \verb'? :'}
        The {\bf? :} ternary infix operator selects one of two values. The first
        argument must be type "log". The second and third arguments may be any type
        but must be the same type. If the test expression is immediate mode the
        availability of the result is the availability of the selected (2nd or 3rd)
        operand. If the test expression is target mode the availability of the
        result is the availability of the test expression ANDed with the logical
        result of the test expression and the availability of the selected 2nd or
        3rd operand. Thus as long as the test expression is available and the
        selected 2nd or 3rd operand is available the result will be available even
        if the unselected operand is not available. This operator may be nested.

        \index{operator!binary \verb'>>'}
        \index{operator!binary \verb'<<'}
        The {\bf\verb'>>'} and {\bf\verb'<<'} binary infix shift operators can
        only be used on immediate types "uint" or "int" or target types "bits:",
        "uint:", "int:", "fixed:" or "ufixed:". The second argument may be an
        immediate variable of type "uint" or "int" (but cannot have a negative
        value), or a target variable of type "uint:". Right shifts on signed
        values extend the sign bit.

        \index{operator!binary \verb'>>^'}
        \index{operator!binary \verb'^<<'}
        The {\bf\verb'>>^'} and {\bf\verb'^<<'} binary infix rotate operators
        can only be used on target mode values which can be of any type. The
        second argument must be an immediate variable of type "uint" or "int"
        and cannot have a negative value. These operators perform a simple
        bit rotation and take no account of the data type.
        

        \index{operator!assignment}
        The following assignment symbols are not treated as operators, but
        rather as part of statement syntax.
        
        \begin{tabular}{| l | l |}
        \hline
        {\bf =}          & immediate mode assignment \\
        \hline
        {\bf :=}         & value, priority or clock mode assignment \\
        \hline
        {\bf \verb'<-'}  & static mode assignment \\
        \hline
        {\bf \verb'<<-'} & queue mode assignment \\
        \hline
        {\bf \verb'|-'}  & selectvalue or conditional output mode assignment \\
        \hline
        \end{tabular}
        
        The {\bf ++} and {\bf --} postfix
        operators can also be used as a form of assignment (see sections
        \autorefph{sec:im_ass} and \autorefph{sec:stat_ass}).

        The following C operators have not been included (or in the case of
        \verb'->' the semantics are different).

        \index{operator!omitted C}
        \begin{tabular}{| l | l | l |}
        \hline
        \verb',' & binary & sequential evaluation \\
        \hline
        \verb'->' & binary & field of variable pointed to \\
        \hline
        \verb'*' & unary & value of variable pointed to \\
        \hline
        \end{tabular}



    \section{Expressions}

        \index{expression}
        \index{value!expression}
        An expression may contain
        \blist
        \item   constants
        \item   variables of primitive type
        \item   function calls returning primitive type
        \item   operators
        \item   parentheses
        \blend
        A primitive type may be a simple variable or a field or array member
        from a compound nested type which is a primitive type. (i.e. not
        arrays or structures as such). Variables may be {\bf immediate}, {\bf
        value}, {\bf selectvalue}, {\bf static}, {\bf input} or {\bf queue} mode.

        \subsection{expression evaluation}

            \index{expression!evaluation}
            All operators are combinatorial. Binary logical operators \verb'&&' and
            \verb'||' {\bf are} short-circuit evaluated when the first operand is
            immediate. Binary logical operators \verb'&&' and \verb'||' {\bf are
            not} short-circuit evaluated when the first operand is a target mode
            since that would imply sequential evaluation.

            If the first operand of the conditional ternary operator \verb'?:'
            is immediate mode, only one of the following two operands is
            evaluated, i.e. this operator is also short-circuit evaluated.

            An expression is available when all queues referenced in the
            expression are available. Where the ternary operator {\bf \verb'?:'}
            is used, all three operands must be available.

        \subsection{mixed types}

            \index{expression!variable modes}
            The arithmetic operators '+', '-', '*', '/' and '\verb"%"' accept
            the types shown on the axes of the following table, the table
            entries giving the result type. Types shown with an appended ":" are
            target types while those without are immediate.

            \begin{tabular}{| l || l | l | l | l | l | l | l | l |}
            \hline
                    & uint    & int    & float & str & uint:   & int:   & ufixed: & fixed: \\
            \hline \hline
            uint    & uint    & int    & float & str & uint:   & int:   & ufixed: & fixed: \\ \hline
            int     & int     & int    & float & str & int:    & int:   & fixed:  & fixed: \\ \hline
            float   & float   & float  & float & str & -       & -      & -       & -      \\ \hline
            str     & str     & str    & str   & str & -       & -      & -       & -      \\ \hline
            uint:   & uint:   & int:   & -     & -   & uint:   & int:   & ufixed: & fixed: \\ \hline
            int:    & int:    & int:   & -     & -   & int:    & int:   & fixed:  & fixed: \\ \hline
            ufixed: & ufixed: & fixed: & -     & -   & ufixed: & fixed: & ufixed: & fixed: \\ \hline
            fixed:  & fixed:  & fixed: & -     & -   & fixed:  & fixed: & fixed:  & fixed: \\ \hline
            \end{tabular}

            The remainder operator '\verb"%"' only accepts types {\bf int}, {\bf
            uint}, {\bf int:} or {\bf uint:}.

            The addition operator '+' where at least one operand is a string
            concatenates the operands producing a string result. The other
            arithmetic operators do not accept string operands.

            Arithmetic operators for target float types are not
            implemented\footnote{Arithmetic expressions cannot use target
            floating point values because floating point operations are
            sufficiently complicated that they must be pipelined to achieve a
            worthwhile execution rate. By definition expression operators must
            be able to execute in a single clock cycle if the operands are
            available. Single cycle target floating point operators could be
            implemented, but the clock rate at which they were used would have
            to be very low.}. Target floating point operations are carried out
            by pipelined library modules.
            
            Addition and subtraction of target fixed point operands with
            immediate floating point values is allowed, the floating point
            values being converted to fixed point with the same number of
            fractional bits, lesser significant bits being rounded off.

            Logical operators only accept {\em log} operands. The result type is
            also {\em log}.

            Relational operators accept {\em int}, {\em uint} and {\em str}
            operands. The \verb'==' and \verb'!=' operators will also accept
            {\em log}, {\em bits} and {\em type} operands and can also compare
            clock mode variables. The result is type {\em log}. Relational
            operators with arguments of type {\em str} perform a lexical
            comparison.

            Shift operators accept {\em bits}, {\em int}, {\em uint} {\em
            fixed} and {\em ufixed} for the left operand but only {\em int} and
            {\em uint} for the right operand (shift count). The result type is
            the type of the left operand.

            Bitwise operators accept {\em fixed}, {\em ufixed}, {\em int}, {\em
            uint}, {\em log} and {\em bits} for the operands. Where the
            operands are of different size the smaller will be expanded to match
            the larger, extending the sign bit in the case of signed types. The
            result width will be the same as that of the larger operand. The
            result type is determined by the following in order -
            
            \blist
            \item   If either operator is {\em fixed} then the other operator
                    cannot be {\em fixed} or {\em ufixed} and the result type
                    will be {\em fixed}.                    
            \item   If either operator is {\em ufixed} then the other operator
                    cannot be {\em fixed} or {\em ufixed} and the result type
                    will be {\em ufixed}.                    
            \item   If either operand is {\em int} the result is {\em int}.
            \item   if either operand is {\em uint} the result is {\em uint}.
            \item   if either operand is {\em bits} the result is {\em bits}.
            \blend

        \subsection{expression width}

            \index{expression!width}
            Immediate integer variables and expressions have a 64 bit value.
            Immediate floating point variables and expressions have a 52.11
            format (IEEE 754 double precision). Target integer and floating
            point values may be any size.
        
            Where signed and unsigned operands are mixed the unsigned operand is
            zero padded by one bit to render it signed.

            For mixed immediate and target modes the immediate operand will be
            converted to a target operand. The width will be the minimum which
            can contain the value, assuming a sign bit in the case of {\em int}.

            For addition and subtraction the result width is the widest operand
            width plus one. For fixed point types the result point offset is
            equal to the largest or only operand point offset.

            For multiplication the result width is the sum of the operand
            widths. For fixed point types the result point offset is the sum of
            the operand point offsets (integer types have zero point offset).

            For division the result width is the sum of the operand widths. For
            fixed point types the result point offset is the first operand point
            offset plus the second operand width minus the second operand point
            offset (see example below).

            For remainder the result is the same width and type as the second
            operand.

            For '\verb'>>'' the result width is the same as the first operand
            width. For '\verb'<<'' the result width is the first operand width
            plus the shift count (immediate) or maximum possible shift count
            (target). For bitwise logical operators the result width is the
            maximum of the operand widths. The result type is the same as that
            of the left operand.

            For '?:' the result width is the width of the widest of the other
            two operands.

            For bitwise logical operators the operands may be any width, the the
            smaller being padded to the size of the larger (with sign extension
            if int). The result is the bitwise AND, OR or XOR of each matching
            pair of bits, one from each operand.
            
            Fixed point divider example -
            \begin{lstlisting}[gobble=12, language=threepl]
               static(num, "fixed:8.5");
               static(den, "fixed:10.4");
               static(quot, "fixed:12.15");
               quot <- num / den;
            \end{lstlisting}
            Note: this code really should use divide(quot <- num, den)\footnote{
            While division has been included in the set of arithmetic
            operations, current FPGAs do not provide division hardware elements.
            As a consequence 3PL implements division as a cascaded series of
            adder/subtracters resulting in a very great logic depth. For data to
            propagate through such a divider in one clock cycle would imply a
            very low clock rate. To avoid this problem stdlib.3pl includes a
            pipelined divide library divlib.3pl whih contains a pipelined divide
            module. This module uses the same logic as produced for the divide
            operator but each stage is pipelined. The compiler itself cannot
            generate such a module implicitly since the language requires basic
            operations to be executed in one clock cycle. } 
            The numerator width is 13 and the denominator width 14 so the result
            width will be 27. The result fixed point offset will be $5 + 14 - 4$
            which is 15. Thus the quotient has been given the type
            "fixed:12.15". The type width and offset could be different to the
            division result type in which case the result would be truncated or
            padded/sign-extended as required.
            
            For mixed types -
            \begin{lstlisting}[gobble=12, language=threepl]
               static(num, "ufixed:6.3");
               static(den, "int:3");
               static(quot, "fixed:7.6");
               quot <- num / den;
            \end{lstlisting}
            The numerator width is 9 and the denominator width is 3. Because the
            denominator is signed the numerator will be converted to
            "fixed:7.3", making room for the extra sign bit. The result width
            will be the sum of these which is 13. The result fixed point offset
            will be $3 + 3 - 0$ which is 6. Thus the result type is "fixed:7.6".
            
        
        \subsection{operand clock domains}

            \index{expression!clock domain}
            If an operator has two or three operands the operands must either
            have the same clock domain or have no clock domain.
            Immediate mode variables and expressions have no clock domain
            and value variables initialised to a constant have no clock domain.


        \subsection{assignment}\label{sec:assignments}

            \index{expression!assignment}
            The following table shows allowed right-hand-side types for each
            left-hand-side type. Target types have an appended ':'.

            \begin{tabular}{| l | l r l l l l l l l l l |}
            \hline
            uint    & uint & +int &       &     &      &       &       &      &         &        &        \\ \hline
            int     & uint & int  &       &     &      &       &       &      &         &        &        \\ \hline
            float   & uint & int  & float &     &      &       &       &      &         &        &        \\ \hline
            str     &      &      &       & str & type &       &       &      &         &        &        \\ \hline
            type    &      &      &       & str & type &       &       &      &         &        &        \\ \hline
            bits:   &      &      &       &     &      & bits: &       &      &         &        &        \\ \hline
            uint:   & uint & +int &       &     &      &       & uint: &      &         &        &        \\ \hline
            int:    & uint & int  &       &     &      &       & uint: & int: &         &        &        \\ \hline
            ufixed: & uint & +int & float &     &      &       & uint: &      & ufixed: &        &        \\ \hline
            fixed:  & uint & int  & float &     &      &       & uint: & int: & ufixed: & fixed: &        \\ \hline
            float:  & uint & int  & float &     &      &       &       &      &         &        & float: \\ \hline
            \end{tabular}

            "+int" indicates that the immediate integer value must not be
            negative.
             
            Fixed point types are truncated at the binary point when assigned
            to integer types.
            
            The LHS and RHS widths, point offsets and floating point mantissa
            and biased exponent sizes may differ. Where the LHS type width is
            smaller than that of the RHS the most significant bits are
            truncated. If the LHS width is greater than that of the RHS the
            value will be sign extended or zero padded at the most significant
            end as appropriate.
            
            Note that types "str" and "type" can be assigned to each other.
            Assigning a type to a string variable produces the same string value
            as calling the nested inbuilt functions {\bf
            typestring(typeval(t))}, i.e. the type is converted to the type
            specification string from which it can be generated. Assigning a
            string to a type variable is equivalent to calling inbuilt function
            {\bf strtotype()} \autorefph{sec:ifstot}.
            
            target type "bits:" can only be assigned type "bits:", however the
            {\bf cast()} inbuilt function can be used to cast any other type to "bits:".
            
            The types allowed for assignment generally apply also to variable
            initialisation.

        \subsection{expression optimisation}

            \index{expression!optimisation}
            The following special cases with mixed target mode and constant
            operands are handled -
            \blist
            \item   multiplication by 0 - value is constant 0
            \item   multiplication by 1 - value is the other operand
            \item   multiplication by a power of 2 - use a fixed left shift
                    for a power greater than zero or a fixed right shift
                    for a power less than zero
            \item   division by 0 - a fatal error
            \item   division by 1 - value is the other operand
            \item   division of 0 - value is constant 0
            \item   division by a power of 2 - use a fixed right shift for a
                    power greater than zero or a fixed left shift for a power
                    less than zero
            \item   remainder by 0 - a fatal error
            \item   remainder by 1 - value is constant 0
            \item   remainder of 0 - value is constant 0
            \item   remainder by a power of 2 - use an AND with a mask
            \item   addition of 0 - value is the other operand
            \item   subtraction of 0 - value is the other operand
            \item   subtraction from 0 - converted to a unary - of the other
                    operand
            \item   logical AND of false - value is constant false
            \item   logical AND of true - value is the other operand
            \item   logical OR of false - value is the other operand
            \item   logical OR of true - value is constant true
            \item   logical XOR of false - value is the other operand
            \item   logical XOR of true - converted to a logical NOT of the
                    other operand
            \item   int \verb'<' 0 - value determined from the sign bit only
            \item   int \verb'>= 0' - value determined from the sign bit only
            \item   0 \verb'<=' int - value determined from the sign bit only
            \item   0 \verb'>' int - value determined from the sign bit only
            \blend
            Bitwise operators with mixed target mode and constant operands are
            optimised during netlist generation.
            
            3PL does no other expression optimisation. Target
            expressions should be written the way you want them implemented.
